\SourcePage{367}
\SGAsection{5}{以幺幂群为核的乘法型群扩张}
\textbf{5.1. \Term{theorem}{定理}的陈述。}
\begin{SGAresult}{定义}{5.1.0}
设 $k$ 为域，$G$ 为 $k$-\Term{algebraic-group}{代数群}。按照 Rosenlicht 在《Questions of rationality for solvable algebraic groups over non perfect fields》，Annali di Math. 61 (1963) 中引入的术语，若 $G$ 满足下列\Term{equivalence}{等价}条件，就称其为“$k$-\Term{solvable}{可解}的（$k$-résoluble）”：
\begin{enumerate}
\item[i)] $G$ 有一条\Term{composition-series}{合成列}，其逐次商均\Term{isomorphism}{同构}于 $\mathbf G_a$。
\item[ii)] $G$ 有一条\Term{characteristic-composition-series}{特征合成列} $(G_i)$，使逐次商 $G_i/G_{i+1}$ 交换，且都有一条逐次商\Term{isomorphism}{同构}于 $\mathbf G_a$ 的\Term{composition-series}{合成列}。
\end{enumerate}
\end{SGAresult}
上述文献证明了 i)$\Rightarrow$ii)。实际上，可取 3.9 ii) 的证明中引入的\Term{composition-series}{合成列}作为 $(G_i)$。
\begin{SGAresult}{\Term{theorem}{定理}}{5.1.1}
设 $k$ 为域，$U$ 为\Term{unipotent}{幺幂} $k$-\Term{algebraic-group}{代数群}，$H$ 为\Term{multiplicative-type}{乘法型} $k$-群，$E$ 为以 $U$ 为核的 $H$ 的 $k$-\Term{algebraic-group}{代数群}\Term{extension}{扩张}，即有\Term{exact-sequence}{正合列}
\[1\longrightarrow U\longrightarrow E\longrightarrow H\longrightarrow1.\]
则：
\begin{enumerate}
\item[i)] 在下列各情形中，\Term{extension}{扩张} $E$ 都平凡：
\begin{enumerate}
\item[a)] $k$ \Term{algebraically-closed}{代数闭}；
\item[b)] $k$ \Term{perfect-field}{完美}，且 $U$、$H$ 至少有一个\Term{connected}{连通}；
\item[c)] $U$ 是 $k$-\Term{solvable}{可解}的；
\item[d)] $U$ \Term{smooth}{光滑}且 $H$ \Term{connected}{连通}。
\end{enumerate}
\item[ii)] 若 $H'$ 和 $H''$ 是 $H$ 在 $E$ 中的两个\Term{lifting}{提升}，则在下列各情形中，它们均由 $U(k)$ 中某个元素\Term{conjugacy}{共轭}：
\begin{enumerate}
\item[a)] $k$ \Term{algebraically-closed}{代数闭}且 $H$ \Term{smooth}{光滑}；
\item[b)] $k$ \Term{algebraically-closed}{代数闭}且 $U$ \Term{smooth}{光滑}。
\end{enumerate}
\end{enumerate}
\end{SGAresult}
（在证明过程中，我们还会指出另一些无需假定 $k$ \Term{algebraically-closed}{代数闭}而 ii) 的结论仍成立的情形。）
若 $U$ 为定义在域 $k$ 上的\Term{algebraic-group}{代数群}（相应地，交换\Term{algebraic-group}{代数群}），以 $\mathrm H^1(k,U)$（相应地，$\mathrm H^i(k,U)$，$i\geq0$）表示 $k$ 取值于 $U$ 的 Galois \Term{cohomology}{上同调}带点集（相应地，第 $i$ 个群）（参见 J.-P. Serre，《Cohomologie Galoisienne》，Lecture Notes Maths. 第 5 号，Springer）。
\SourcePage{368}
若 $S$ 为\Term{prescheme}{预概形}，$H$ 为交换 $S$-\Term{group-prescheme}{群预概形}，$G$ 为\Term{action}{作用}于 $H$ 的 $S$-\Term{group-prescheme}{群预概形}，则以 $\mathrm H^i(G,H)$\footnote{编者注：为与附录 I 和第 I 讲的记号一致，已将 $\mathrm H^i_0$ 改为 $\mathrm H^i$。}表示 $G$ 取值于 $H$ 的第 $i$ 个 Hochschild \Term{cohomology-group}{上同调群}（附录 I 1）。
我们将分几步证明 5.1.1。

\textbf{5.2. 在 $U$ \Term{smooth}{光滑}且 $H$ \Term{etale}{平展}的情形下证明 5.1.1 i) 和 ii)。}
\begin{SGAresult}{\Term{lemma}{引理}}{5.2.1}
沿用 5.1.1 的记号。若 $H$ \Term{etale}{平展}，则在下列情形中，\Term{canonical}{典范}\Term{morphism}{态射} $E\to H$ 有一个定义在 $k$ 上的\Term{section}{截面} $s:H\to E$：
a) $k$ \Term{algebraically-closed}{代数闭}；b) $k$ \Term{perfect-field}{完美}，且 $U$ \Term{smooth}{光滑}、\Term{connected}{连通}；c) $U$ 是“$k$-\Term{solvable}{可解}的”。
\end{SGAresult}
\begin{proof}
a) 直接来自 $E(k)\to H(k)$ \Term{surjectivity}{满射}。由 4.1.2 b) 和 2.3（bis），b)$\Rightarrow$c)，故只须处理 c)。设 $x$ 为 $H$ 的一点，则 $x$ 是 $H$ 中既开又闭的部分，诱导\Term{subscheme}{子概形}\Term{isomorphism}{同构}于 $\operatorname{Spec}K$，其中 $K$ 是 $k$ 的\Term{finite}{有限}\Term{separable}{可分}\Term{extension}{扩张}。设 $X$ 为\Term{scheme}{概形} $x$ 在 $E$ 中的\Term{inverse-image}{逆像}。$K$-\Term{scheme}{概形} $X$ 是群 $U_K$ 下的 $K$-\Term{principal-homogeneous-space}{主齐性空间}。$U_K$ 有一条逐次商\Term{isomorphism}{同构}于 $(\mathbf G_a)_K$ 的\Term{composition-series}{合成列}，故 $\mathrm H^1(K,U_K)=0$（J.-P. Serre，《Cohomologie Galoisienne》，第 III 章，命题 6），从而 $X$ 平凡，因而有 $K$-\Term{rational-point}{有理点}。这样，在 $H$ 的每一点 $x$ 上都得到 $E\to H$ 的一个 $k$-\Term{section}{截面}，故存在\Term{section}{截面} $H\to E$。
\end{proof}
\begin{SGAresult}{\Term{lemma}{引理}}{5.2.3}
\footnote{编者注：没有编号 5.2.2。}沿用 5.1.1 的记号，假定 $H$ \Term{etale}{平展}。则：
\begin{enumerate}
\item[i)] 在下列各情形中，\Term{extension}{扩张} $E$ 都平凡：
a) $U$ 交换，且 $E\to H$ 有\Term{section}{截面}；b) $k$ \Term{algebraically-closed}{代数闭}；c) $k$ \Term{perfect-field}{完美}且 $U$ \Term{connected}{连通}；d) $U$ 是“$k$-\Term{solvable}{可解}的”。此外，在上述各情形中，$H$ 在 $E$ 中的任意两个\Term{lifting}{提升} $H'$、$H''$ 均由 $U(k)$ 中某个元素\Term{conjugacy}{共轭}。
\item[ii)] 设 $R:(\mathrm{Sch}/k)^\circ\to\mathrm{Ens}$ 为如下 $k$-\Term{functor}{函子}：对每个 $k$-\Term{prescheme}{预概形} $S$，$R(S)$ 是 $H_S$ 在 $E_S$ 中的\Term{lifting}{提升}的\Term{set}{集合}。若 $U$ 交换，则 $R$ 由\Term{nonempty}{非空}\Term{affine}{仿射} $k$-\Term{scheme}{概形}表示。群 $U$ 以\Term{inner-automorphism}{内自同构}\Term{action}{作用}于 $R$，此\Term{action}{作用}使 $R$ 成为 $U$ 下的\Term{homogeneous-space}{齐性空间}（对于 fpqc \Term{topology}{拓扑}而言，参见第 IV 讲）。
\end{enumerate}
\end{SGAresult}
\begin{proof}[i) 的证明]
先把 b)、c)、d) 化归为 a)。
情形 b)：$U$ 有一条逐次商交换的\Term{characteristic-composition-series}{特征合成列}（3.5 i)$\Leftrightarrow$iv)），故立即化归到 $U$ 交换的情形。
\SourcePage{369}
又由 5.2.1 a)，$E\to H$ 有\Term{section}{截面}，于是归结为 a)。
情形 d)：同样处理，使用 5.1.0 ii) 和 5.2.1 c)。
情形 c)：设 $E_{\mathrm{red}}$ 为 $E$ 的关联\Term{reduced}{既约}\Term{subgroup}{子群}。因 $H$ \Term{smooth}{光滑}，$E_{\mathrm{red}}$ 是以 $U'=U\cap E_{\mathrm{red}}$ 为核的 $H$ 的\Term{extension}{扩张}，可以用 $E_{\mathrm{red}}$ 代替 $E$。但 $U$ \Term{connected}{连通}，故 $U'$ 也\Term{connected}{连通}，且显然是 $E_{\mathrm{red}}$ 的\Term{identity-component}{单位连通分支}，因而\Term{smooth}{光滑}。于是 $U'$ \Term{smooth}{光滑}、\Term{connected}{连通}，$k$ \Term{perfect-field}{完美}，故 $U'$ 是 $k$-\Term{solvable}{可解}的（4.1.4 b)），又归结为 d)。
上述化简表明，在 b)、c)、d) 中可以假定 $U$ 有一条\Term{characteristic-composition-series}{特征合成列} $(U_i)$，使 $U_i/U_{i+1}$ 交换，且映射 $U_i(k)\to(U_i/U_{i+1})(k)$ \Term{surjectivity}{满射}（在 c)、d) 中，后一性质来自 $\mathrm H^1(k,\mathbf G_a)=0$）。直接逐层剖分便知，为证明两个\Term{lifting}{提升} $H'$、$H''$ 的\Term{conjugacy}{共轭}性，只须考虑 $U$ 交换的情形。总之，只须证明 i) a)。这时若 $\mathrm H^2(H,U)=0$，\Term{extension}{扩张} $E$ 就平凡（附录 I 3.1）；若 $\mathrm H^1(H,U)=0$，$H'$ 和 $H''$ 就\Term{conjugacy}{共轭}。为此有下述\Term{lemma}{引理}。
\end{proof}
\begin{SGAresult}{\Term{lemma}{引理}}{5.2.4}
设 $S$ 为\Term{prescheme}{预概形}，$U$ 为交换 $S$-\Term{group-prescheme}{群预概形}，$H$ 为\Term{action}{作用}于 $U$ 的\Term{etale}{平展}、\Term{finite}{有限}、秩为 $n$ 的 $S$-\Term{group-prescheme}{群预概形}。在下列两种情形中，$n$ 均\Term{annihilate}{零化}各群 $\mathrm H^i(H,U)$（$i>0$）：
a) $H$ 为常值 $S$-群（第 I 讲 4.1）；b) $S$ 为域的谱。
\end{SGAresult}
\begin{proof}[a) 的证明]
按假设，$H$ 是与一个阶为 $n$ 的普通群 $\{H\}$ 关联的\Term{constant-group}{常值群}。此时显然\footnote{编者注：这是 M. Demazure 和 P. Gabriel，《Groupes algébriques I》，Masson \& North-Holland (1970) 的命题 III.6.4.2。}$\mathrm H^i(H,U)$ \Term{isomorphism}{同构}于群 $\{H\}$ 取值于普通群 $U(S)$ 的第 $i$ 个\Term{cohomology-group}{上同调群} $\mathrm H^i(\{H\},U(S))$；这些群被 $n$ \Term{annihilate}{零化}是经典结果（J.-P. Serre，《Corps locaux》，第 VIII 章，命题 4，推论 1）。
\end{proof}
\begin{proof}[b) 的证明]
取 $x\in\mathrm H^i(H,U)$（$i>0$），用 $i$-\Term{cocycle}{余循环} $f:H^{(i)}\to U$ 表示它（$H^{(i)}$ 表示 $i$ 个 $H$ 在 $k$ 上的乘积）。若 $K/k$ 为使 $H$ \Term{splitting}{分裂}的\Term{finite}{有限} Galois \Term{extension}{扩张}，由 a)，$nf_K$ 为\Term{coboundary}{余边界}。更精确地，若按下式定义 $K$-\Term{morphism}{态射} $F_K:H_K^{(i-1)}\to U_K$：
\[F_K(h_1,\ldots,h_{i-1})=\sum_{h\in H(K)}f_K(h_1,\ldots,h_{i-1},h),\]
则简单计算给出，在相差一个\Term{signature}{符号}的意义下，
\[d(F_K)=nf_K\qquad(d\text{ 为余边界算子}).\]
直接的 Galois 论证说明，$F_K$ 来自 $k$-\Term{morphism}{态射} $F:H^{(i-1)}\to U$，因而 $d(F)=nf$。
\end{proof}
\SourcePage{370}
\begin{SGAresult}{推论}{5.2.4 bis}
沿用 5.2.4 的记号，再假定 $U$ 在 $S$ 上\Term{flat}{平坦}、\Term{finitely-presented}{有限表现}，且各\Term{fiber}{纤维}\Term{unipotent}{幺幂}，$n$ 与 $S$ 的所有\Term{residue-characteristic}{剩余特征}互素。则在上述 a)、b) 中，对 $i>0$ 都有 $\mathrm H^i(H,U)=0$。
\end{SGAresult}
\begin{proof}
只须证明 $U$ 中的 $n$ 次幂映射为\Term{isomorphism}{同构}，因为这将使 $\mathrm H^i(H,U)$ 上的乘 $n$ 同时是\Term{isomorphism}{同构}和零\Term{morphism}{态射}，故该群为零。按 $U$ 的假设，只须在\Term{fiber}{纤维}上检查 $n$ 次幂映射为\Term{isomorphism}{同构}（EGA IV 17.9.5），因而可设 $S$ 为\Term{characteristic}{特征} $p$ 的域 $k$ 的谱。因 $(n,p)=1$，$U$ 中的 $n$ 次幂映射为\Term{etale}{平展}\Term{morphism}{态射}（第 VII 讲），又为\Term{monomorphism}{单态射}（2.4 i)），故为\Term{open-immersion}{开浸入}（EGA IV 17.9.1）。这已证明它在\Term{identity-component}{单位连通分支} $U^0$ 上的\Term{restriction}{限制}为\Term{isomorphism}{同构}；再因 $U/U^0$ 为\Term{finite}{有限} $p$-群，结论得证。
\end{proof}
这也完成 5.2.3 i) a) 的证明，因为 $H$ 为\Term{etale}{平展}\Term{multiplicative-type}{乘法型}群，故其阶与 $p$ 互素。
\begin{proof}[5.2.3 ii) 的证明]
显然 $R$ 是 fpqc \Term{topology}{拓扑}下的\Term{sheaf}{层}。又由\Term{affine}{仿射}\Term{scheme}{概形}的\Term{descent}{下降}，5.2.3 ii) 的断言是 fpqc \Term{local-scheme}{局部}的，故可以假定 $k$ \Term{algebraically-closed}{代数闭}。这时 $H$ \Term{splitting}{分裂}，\Term{extension}{扩张} $E$ 平凡（i) b)）。取 $H$ 在 $E$ 中的一个\Term{lifting}{提升} $H'$。对每个 $k$-\Term{prescheme}{预概形} $S$ 及 $H_S$ 在 $E_S$ 中的每个\Term{lifting}{提升} $H''$，因 $\mathrm H^1(H_S,U_S)=0$（5.2.4 bis），$H'_S$ 和 $H''$ 由 $U(S)$ 中某元素\Term{conjugacy}{共轭}。设 $U^{H'}$ 为 $U$ 在 $H'$ 下的\Term{invariants}{不变量}\Term{sheaf}{层}，它由 $U$ 的一个\Term{algebraic-subgroup}{代数子群}表示（第 VIII 讲 6.5 d)）。由上述讨论，$k$-\Term{morphism}{态射}
\[U\longrightarrow R,\qquad u\longmapsto\operatorname{int}(u)H'\quad(u\in U(S))\]
定义一个 $k$-\Term{isomorphism}{同构} $U/U^{H'}\xrightarrow{\sim}R$。这证明了 $R$ 的\Term{representability}{可表示性}及其\Term{affine}{仿射}性（2.1）。
\end{proof}
\textbf{注记 5.2.5。}可以证明，即使 $U$ 不交换，5.2.3 ii) 的断言仍成立，但证明 5.1.1 无需用到这一点。

\textbf{5.3. $H$ \Term{smooth}{光滑}的情形。}
\begin{SGAresult}{命题}{5.3.1}
将假设“$H$ \Term{etale}{平展}”改为“$H$ \Term{smooth}{光滑}”后，5.2.3 i) 的各项断言仍成立。
\end{SGAresult}
\begin{proof}
与 5.2.3 i) 的证明一样，可进一步假定 $U$ 交换。设 $N'$ 为与 $p$ 互素的正整数集，按整除关系排序。对每个 $n\in N'$，${}_nH'$ 为\Term{etale}{平展}群，而族 $({}_nH')_{n\in N'}$ 在 $H$ 中\Term{schematically-dense}{概形稠密}，因为 $H$ \Term{smooth}{光滑}（第 IX 讲 4.10）。以 $E_n$ 表示 ${}_nH$ 在 $E$ 中的\Term{inverse-image}{逆像}，则 $E_n$ 是以 $U$ 为核的 ${}_nH$ 的\Term{extension}{扩张}。最后设 $R_n$ 为 ${}_nH$ 在 $E_n$ 中的\Term{lifting}{提升}的 $k$-\Term{functor}{函子}（参见 5.2.3 ii)）。若 $n\mid m$，显然有自然 $k$-\Term{morphism}{态射} $R_m\to R_n$，故各 $R_n$ 构成 $k$-\Term{functor}{函子}的\Term{projective-system}{投射系}。
\SourcePage{371}
因 $R_n$ 由\Term{nonempty}{非空}\Term{affine}{仿射} $k$-\Term{scheme}{概形}表示（5.2.3 ii)），且非零环的\Term{filtered}{滤过}\Term{inductive-limit}{归纳极限}非零，\Term{functor}{函子} $R=\varprojlim R_n$ 也由\Term{nonempty}{非空}\Term{affine}{仿射} $k$-\Term{scheme}{概形}表示（EGA IV 8 及 1.9.1）。故存在 $k$ 的\Term{extension}{扩张} $K$ 和一点 $u\in R(K)$。$u$ 在 $R_n(K)$ 中的像 $u_n$ 对应于 $({}_nH)_K$ 在 $(E_n)_K$ 中的\Term{lifting}{提升} $H'_n$。由构造，当 $n\mid m$ 时 $H'_n={}_n(H'_m)$。令 $U_n=(U_K)^{H'_n}$。选取 $H'_n$ 便可将 $(R_n)_K$ 与 $U_K/U_n$ 等同。但 $H'_n$ 构成递增\Term{filtered-family}{滤过族}，故 $U_n$ 构成递减\Term{filtered-family}{滤过族}，并因此在 $n$ \Term{large-relative-to-cardinal}{足够大}时\Term{stable}{稳定}（$U_K$ 是\Term{noetherian}{诺特}的）。于是 $(R_n)_K$ 族\Term{stable}{稳定}，故 $R_n$ 族也\Term{stable}{稳定}。总之，对充分大的 $n$，$R_n=R$。
按现有假设，由 5.2.3 i)，$R_n(k)$ \Term{nonempty}{非空}，故可找到 ${}_nH$（$n\in N'$）的一族相容\Term{lifting}{提升} $H'_n$。设 $H'$ 为 $E$ 中包含所有 $H'_n$ 的最小闭\Term{subscheme}{子概形}（第 VI\textsubscript{B} 讲 §7）。第 XV 讲 4.6 的论证表明，$H'$ 为\Term{smooth}{光滑}交换\Term{algebraic-group}{代数群}，且其形成与\Term{base-field}{基域}\Term{extension}{扩张}交换。因此，为证明 $H'$ 为 $H$ 在 $E$ 中的\Term{lifting}{提升}，可以假定 $k$ \Term{algebraically-closed}{代数闭}。由 BIBLE 4，\Term{theorem}{定理} 4，$H'$ 是一个（\Term{smooth}{光滑}）\Term{multiplicative-type}{乘法型}群 $M$ 与\Term{unipotent}{幺幂}群 $V$ 的\Term{direct-product}{直积}。各 $H'_n$ 必然包含在 $M$ 中（2.4），由 $H'$ 的定义便有 $H'=M$。故 $H'$ 为\Term{multiplicative-type}{乘法型}，从而 $H'\cap U=0$。因此\Term{morphism}{态射} $H'\to H$ 是\Term{monomorphism}{单态射}，又由\Term{dense}{稠密}性\Term{theorem}{定理}（第 IX 讲 4.10）为\Term{epimorphism}{满态射}，故确为\Term{isomorphism}{同构}。
现设 $H'$、$H''$ 为 $H$ 在 $E$ 中的两个\Term{lifting}{提升}。对每个 $n\in N'$，以 $T_n=\operatorname{Transp}({}_nH',{}_nH'')$ 表示从 ${}_nH'$ 到 ${}_nH''$ 的\Term{transporter}{输运子}，它由 $U$ 的闭\Term{subscheme}{子概形}表示（第 VIII 讲 6.5 e)）。各 $T_n$ 构成递减\Term{filtered-family}{滤过族}，由 5.2.3 i) a) 均\Term{nonempty}{非空}。设 $T$ 为其\Term{stable}{稳定}值。在 5.2.3 i) 的假设下，$T_n(k)$ \Term{nonempty}{非空}，故存在 $u\in U(k)$ 使所有 $n\in N'$ 都满足 ${}_nH''=\operatorname{int}(u){}_nH'$。于是 $H''=\operatorname{int}(u)H'$（第 IX 讲 4.8 b)）。
\end{proof}

\textbf{5.4. $U$ \Term{radicial}{根性}的情形。}
\begin{SGAresult}{命题}{5.4.1}
若 $k$ 为\Term{characteristic}{特征} $p>0$ 的\Term{perfect-field}{完美}域，且 $U$ \Term{radicial}{根性}，则 5.1.1 中的\Term{extension}{扩张} $E$ 平凡。
\end{SGAresult}
\begin{proof}
利用 $U$ 的\Term{characteristic-composition-series}{特征合成列}，可限于 $U=(\alpha_p)^r$ 的情形（3.9）。由附录 II 2.2 和 2.1，有 $k$-\Term{functor}{函子}的\Term{isomorphism}{同构}
\[\operatorname{Aut}_{k\text{-gr}}(\alpha_p)^r\xrightarrow{\sim}
\operatorname{Aut}_{p\text{-Lie}}(\operatorname{Lie}(\alpha_p)^r)\xrightarrow{\sim}
\operatorname{GL}(\operatorname{Lie}(\alpha_p)^r).\]
取秩为 $r$ 的 $k$-\Term{vector-space}{向量空间} $V$ 及向量 $k$-\Term{group-scheme}{群概形} $\mathbf W(V)$（第 I 讲 4.6），将 $(\alpha_p)^r$ 与 ${}_FV$ 等同。上述讨论表明，由\Term{extension}{扩张} $E$ 定义的 $H$ 在 ${}_FV$ 上的\Term{action}{作用}来自 $H$ 在 $V$ 中的\Term{linear-representation}{线性表示} $v$。考虑\Term{exact-sequence}{正合列}
\[0\longrightarrow{}_FV\longrightarrow V\xrightarrow{F}V^{(p)}\longrightarrow0.\tag{*}\]
\SourcePage{372}
其中 $F$ 为 Frobenius \Term{morphism}{态射}。只要用\Term{linear-representation}{线性表示} $H\xrightarrow v\operatorname{GL}(V)\xrightarrow F\operatorname{GL}(V^{(p)})$ 定义 $H$ 在 $V^{(p)}$ 上的\Term{action}{作用}，(*) 就是 $H$-群的\Term{exact-sequence}{正合列}。
因 $k$ \Term{perfect-field}{完美}，$F:V\to V^{(p)}$ 诱导\Term{surjectivity}{满射} $V(k)\to V^{(p)}(k)$。由附录 I 2.1，(*) 定义\Term{exact-sequence}{正合列}
\[\mathrm H^1(H,V^{(p)})\longrightarrow\operatorname{Ext}_{\mathrm{alg}}(H,{}_FV)\longrightarrow\operatorname{Ext}_{\mathrm{alg}}(H,V).\tag{**}\]
证明 $\operatorname{Ext}_{\mathrm{alg}}(H,V)=0$。取\Term{algebraic-group}{代数群}\Term{extension}{扩张}
\[1\longrightarrow V\longrightarrow E'\xrightarrow h H\longrightarrow1.\]
\Term{scheme}{概形} $E'$ 是底为 $H$、群为 $\mathbf G_a^r$ 的\Term{torsor}{挠子}，故定义 $\mathrm H^1(H,\mathscr O_H^r)$ 的一个元素（这里指\Term{coherent-sheaf}{凝聚层}\Term{cohomology}{上同调}）。因 $H$ \Term{affine}{仿射}，$\mathrm H^1(H,\mathscr O_H^r)=0$（EGA III §1），即 $E'\to H$ 有\Term{section}{截面}。因此 $\operatorname{Ext}_{\mathrm{alg}}(H,V)\simeq\mathrm H^2(H,V)$（附录 I 3.1）。而当 $i>0$ 时，$\mathrm H^i(H,V)=\mathrm H^i(H,\mathbf W(V))=0$（第 IX 讲 3.1）。由 (**) 即得 $\operatorname{Ext}_{\mathrm{alg}}(H,{}_FV)=0$，故 $E$ 是平凡\Term{extension}{扩张}。
\end{proof}

\textbf{5.5. 5.1.1 i) 的证明。}若 $k$ 的\Term{characteristic}{特征}为零，则 $U$ 是 $k$-\Term{solvable}{可解}的（4.1.3），且 $H$ \Term{smooth}{光滑}；由 5.3.1 和 5.2.3 d)，$E$ 为平凡\Term{extension}{扩张}。同样可证 $H$ 在 $E$ 中的两个\Term{lifting}{提升}由 $U(k)$ 中某元素\Term{conjugacy}{共轭}。
以下均假定 $k$ 的\Term{characteristic}{特征}为 $p>0$。
\textbf{i) b) 的证明：$k$ \Term{perfect}{完美}、$U$ \Term{connected}{连通}的情形。}
将化归为 $U$ \Term{radicial}{根性}的情形。因 $k$ \Term{perfect-field}{完美}，$H_{\mathrm{red}}$ \Term{smooth}{光滑}；设 $E'$ 为它在 $E$ 中的\Term{inverse-image}{逆像}。由 5.3.1 和 5.2.3 i) c)，\Term{extension}{扩张}
\[1\longrightarrow U\longrightarrow E'\longrightarrow H_{\mathrm{red}}\longrightarrow1\]
平凡。取 $H_{\mathrm{red}}$ 在 $E$ 中的一个\Term{lifting}{提升} $H_1$。由附录 II 3.1，存在整数 $n>0$ 使 $E^{(n)}=E/{}_{F^n}E$ \Term{smooth}{光滑}；设 $E''$ 为 $E$ 中由 $H_1$ 和 ${}_{F^n}E$ 生成的\Term{algebraic-subgroup}{代数子群}（即 $H_1$ 在 $E^{(n)}$ 中的像在 $E$ 中的\Term{inverse-image}{逆像}）。以 $H''$ 表示 $E''$ 在 $H$ 中的像。断言 $H''=H$。事实上，设 $R$ 为 ${}_{F^n}E$ 在 $H$ 中的像，则 $H''$ 由 $R$ 和 $H_{\mathrm{red}}$ 生成。群 $H/R$ 是 $E^{(n)}$ 的商，故\Term{smooth}{光滑}；因此\Term{canonical}{典范}\Term{morphism}{态射} $H_{\mathrm{red}}\to H/R$ 为\Term{epimorphism}{满态射}，故 $H$ 由 $R$ 和 $H_{\mathrm{red}}$ 生成，从而等于 $H''$。得到\Term{exact-sequence}{正合列}
\[1\longrightarrow U''=U\cap E''\longrightarrow E''\longrightarrow H\longrightarrow1.\tag{$\dagger$}\]
$E''$ 的底\Term{space}{空间}与 $H_1$ 相同，故 $U''$ \Term{radicial}{根性}。显然只须证明\Term{extension}{扩张} ($\dagger$) 平凡，这由 5.4.1 得出。
\textbf{i) b) 的证明：$k$ \Term{perfect}{完美}、$H$ \Term{connected}{连通}的情形。}
\SourcePage{373}
群 $U$ 是以\Term{identity-component}{单位连通分支} $U^0$ 为核的一个\Term{etale}{平展}群的\Term{extension}{扩张}。$U$ \Term{connected}{连通}的情形已处理，故只须考察 $U$ \Term{etale}{平展}的情形。此时有以下更精确的\Term{lemma}{引理}。
\begin{SGAresult}{\Term{lemma}{引理}}{5.5.1}
沿用 5.1.1 的记号，再假定 $U$ \Term{etale}{平展}。则：
i) 若 $H$ \Term{connected}{连通}，$H$ 在 $E$ 中有唯一\Term{lifting}{提升}，即 $E$ 的\Term{identity-component}{单位连通分支} $E^0$。
ii) 若 $k$ \Term{algebraically-closed}{代数闭}，$E$ 平凡，且 $H$ 在 $E$ 中的两个\Term{lifting}{提升}由 $U(k)$ 中某个元素\Term{conjugacy}{共轭}。
\end{SGAresult}
\begin{proof}
i) 因\Term{identity-component}{单位连通分支}的形成与\Term{base-field}{基域}\Term{extension}{扩张}交换，可以限于 $k$ \Term{algebraically-closed}{代数闭}的情形。若 $H$ 为\Term{torus}{环面}，则 $E$ 平凡（5.3.1 和 5.2.3 i) b)），显然 $E^0$ 是 $H$ 的唯一\Term{lifting}{提升}。这已证明在一般情形中，$E^0\cap U$ \Term{radicial}{根性}；它又\Term{etale}{平展}（因 $U$ \Term{etale}{平展}），故为零。于是 $E^0\to H$ 为\Term{monomorphism}{单态射}，又\Term{flat}{平坦}（因 $E^0$ 在 $E$ 中开）且\Term{surjectivity}{满射}（因 $H$ \Term{connected}{连通}），故为\Term{isomorphism}{同构}。若 $H'$ 为 $H$ 的另一个\Term{lifting}{提升}，则它\Term{connected}{连通}，故包含在 $E^0$ 中，从而等于 $E^0$。
ii) 设 $H^0$ 为 $H$ 的\Term{identity-component}{单位连通分支}。由 i)，$E^0$ 是 $H^0$ 在 $E$ 中的唯一\Term{lifting}{提升}。令 $E'=E/E^0$、$H'=H/H^0$，有\Term{extension}{扩张}
\[1\longrightarrow U\longrightarrow E'\longrightarrow H'\longrightarrow1.\]
因 $H'$ \Term{etale}{平展}，此\Term{extension}{扩张}平凡（5.2.3 i) b)）。若 $H'_1$ 为 $H'$ 在 $E'$ 中的\Term{lifting}{提升}，$H_1$ 为其在 $E$ 中的\Term{inverse-image}{逆像}，显然 $H_1$ \Term{lifting}{提升} $H$。若 $H_2$ 为 $H$ 在 $E$ 中的另一个\Term{lifting}{提升}，则它包含 $E^0$；由 5.2.3 i) b)，$H_2$ 在 $E'$ 中的像与 $H'_1$ 由某个 $u\in U(k)$ \Term{conjugacy}{共轭}，立即得 $H_2=\operatorname{int}(u)H_1$。
\end{proof}
\textbf{注记 5.5.2。}在 5.5.1 i) 的假设下，容易看出 $E^0$ \Term{centralizes}{中心化} $U$。
\textbf{5.1.1 i) a) 的证明。}利用\Term{composition-series}{合成列}
\[1\longrightarrow U^0\longrightarrow U\longrightarrow U/U^0\longrightarrow1,\]
由 i) b) 和 5.5.1 ii) 一起即得 i) a)。
在证明 5.1.1 c)、d) 之前，先证明 5.1.1 ii)。
\textbf{5.5.1 ii) a) 的证明。}由于尚无满意的一般陈述，以下列出当 $H$ \Term{smooth}{光滑}时 $H$ 在 $E$ 中的两个\Term{lifting}{提升}\Term{conjugacy}{共轭}的若干情形。
\begin{SGAresult}{命题}{5.6.1}
沿用 5.1.1 的记号，另假定 $H$ \Term{smooth}{光滑}，设 $H'$ 和 $H''$ 为 $H$ 在 $E$ 中的两个\Term{lifting}{提升}。则在下列各情形中，它们均由 $U(k)$ 中某个元素\Term{conjugacy}{共轭}：
a) $k$ \Term{algebraically-closed}{代数闭}；b) $U$ 交换；c) $k$ \Term{perfect-field}{完美}且 $U$ \Term{connected}{连通}；
\SourcePage{374}
d) $U$ 是 $k$-\Term{solvable}{可解}的；e) $H'$ 在 $U$ 中的\Term{centralizer}{中心化子}是 $k$-\Term{solvable}{可解}的；f) \Term{multiplicative-type}{乘法型}群 $H$ 被 $k$ 的一个\Term{degree}{次数}与 $p$ 互素的\Term{finite}{有限} Galois \Term{extension}{扩张} $K$ \Term{trivialization}{平凡化}。
\end{SGAresult}
\begin{proof}
a)、b)、c)、d) 来自 5.3.1。
e) 设 $Z$ 为 $H'$ 在 $U$ 中的\Term{centralizer}{中心化子}，$T$ 为从 $H'$ 到 $H''$ 的\Term{transporter}{输运子}。由 a)，$T$ \Term{nonempty}{非空}，故为 $Z$ 下的\Term{principal-homogeneous-space}{主齐性空间}，而对 $Z$ 的假设保证它平凡。
f) 与 5.3.1 一样，只须考虑 $H$ \Term{etale}{平展}的情形。先设 $H$ \Term{diagonalizable}{可对角化}，由阶为 $m$、与 $p$ 互素的普通群 $M$ 定义。给定两个\Term{lifting}{提升} $H'$、$H''$，便确定 $M$ 取值于 $U(k)$ 的一个 $1$-\Term{cocycle}{余循环}，即映射 $h:M\to U(k)$，使对每对 $m,n\in M$，有 $h(mn)=h(m)\,{}^mh(n)$。$H'$、$H''$ 由 $U(k)$ 中某个元素\Term{conjugacy}{共轭}，当且仅当存在 $a\in U(k)$ 使
\[h(m)=a^{-1}({}^ma).\]
抽象群 $U(k)$ 有一条逐次商交换且被 $p$ 的某个\Term{nilpotent}{幂零}化的\Term{composition-series}{合成列}（由 5.6.1 c)，不妨假定 $p>0$）。立即\Term{pushout}{推出}这种情况下 $h$ 是\Term{coboundary}{余边界}。
现考察一般情形。仍以 $Z$ 表示 $H'$ 在 $U$ 中的\Term{centralizer}{中心化子}，$T$ 表示从 $H'$ 到 $H''$ 的\Term{transporter}{输运子}，后者为 $Z$ 下的\Term{torsor}{挠子}。按假设，有 Galois 群为 $G$、\Term{degree}{次数}与 $p$ 互素的\Term{finite}{有限} Galois \Term{extension}{扩张} $K/k$ \Term{trivialization}{平凡化} $H$。由上述讨论，$H'_K$ 和 $H''_K$ 由 $U(K)$ 中某元素\Term{conjugacy}{共轭}，故 $T_K$ 平凡。因此 $T$ 由 $\mathrm H^1(G,Z)$ 中一个元素定义。与上面相同的理由及对 $G$ 的假设给出 $\mathrm H^1(G,Z)=e$，故 $T$ 平凡。
\end{proof}

\textbf{5.7. 5.1.1 ii) b) 的证明。}
\begin{SGAresult}{\Term{lemma}{引理}}{5.7.1}
设 $S$ 为\Term{prescheme}{预概形}，$G$ 为在 $S$ 上\Term{separated}{分离}、\Term{smooth}{光滑}的 $S$-\Term{group-prescheme}{群预概形}，$H$ 为\Term{action}{作用}于 $G$ 的\Term{multiplicative-type}{乘法型} $S$-群。则 $G$ 在 $H$ 下的\Term{invariants}{不变量} $S$-\Term{functor}{函子} $Z=G^H$ 由 $G$ 的一个在 $S$ 上闭且\Term{smooth}{光滑}的 $S$-子\Term{group-prescheme}{群预概形}表示。
\end{SGAresult}
\begin{proof}
$G^H$ 由闭子\Term{group-prescheme}{群预概形}表示，这来自第 VIII 讲 6.5 d)。考虑\Term{semidirect-product}{半直积} $K=G\times_SH$，则 $H$ 在 $K$ 中的\Term{centralizer}{中心化子}等于 $Z\times_SH$。为证明 $Z$ \Term{smooth}{光滑}，须验证：若 $S$ \Term{affine}{仿射}，$S_0$ 为 $S$ 中由平方为零的\Term{ideal}{理想}定义的闭\Term{subscheme}{子概形}，$u_0\in Z(S_0)$，则存在\Term{lifting}{提升} $u_0$ 的 $u\in Z(S)$。因 $G/S$ \Term{smooth}{光滑}，有 $u\in G(S)$ \Term{lifting}{提升} $u_0$。设 $i$ 为 $H$ 到 $K$ 的\Term{canonical}{典范}\Term{immersion}{浸入}，考虑两个 $S$-群\Term{morphism}{态射}
\[H\rightrightarrows K,\qquad i\quad\text{及}\quad\operatorname{int}(u)i=j.\]
因 $u_0\in Z(S_0)$，有 $i_{S_0}=j_{S_0}$。由第 IX 讲 3.2，存在 $v\in K(S)$，\Term{lifting}{提升} $K(S_0)$ 的\Term{unit-section}{单位截面}，且满足
\SourcePage{375}
\[i=\operatorname{int}(v)j=\operatorname{int}(vu)i.\]
故 $vu=(u',v')$ 属于 $(Z\times_SH)(S)$，所以 $u'\in Z(S)$，且\Term{lifting}{提升} $u_0$。
\end{proof}
\begin{SGAresult}{\Term{lemma}{引理}}{5.7.2}
设 $k$ 为域，$H$ 为 $k$-\Term{algebraic-group}{代数群}，$P(H)$ 为下列性质：“对任意\Term{smooth}{光滑}\Term{unipotent}{幺幂} $k$-群 $U$ 及以 $U$ 为核的 $H$ 的任意\Term{extension}{扩张} $E$，$H$ 在 $E$ 中的两个\Term{lifting}{提升}由 $U(k)$ 中某元素\Term{conjugacy}{共轭}。”若 $H$ 是以\Term{multiplicative-type}{乘法型}群 $H'$ 为核的\Term{multiplicative-type}{乘法型}群 $H''$ 的\Term{algebraic-group}{代数群}\Term{extension}{扩张}，且 $P(H')$、$P(H'')$ 成立，则 $P(H)$ 成立。
\end{SGAresult}
\begin{proof}
设 $U$ 为\Term{smooth}{光滑}\Term{unipotent}{幺幂} $k$-群，$E$ 为以 $U$ 为核的 $H$ 的\Term{extension}{扩张}，$H_1$、$H_2$ 为 $H$ 在 $E$ 中的两个\Term{lifting}{提升}，$H'_1$、$H'_2$ 为对应的 $H'$ 的\Term{lifting}{提升}。因 $P(H')$ 成立，有 $u\in U(k)$ 使 $H'_2=\operatorname{int}(u)H'_1$。用 $\operatorname{int}(u)H_1$ 代替 $H_1$ 后，可假定 $H'_1=H'_2$，将其简记为 $H'$。设 $E'=\operatorname{Centr}_E H'$，它等于
\[U^{H'}\cdot H_1=U^{H'}\cdot H_2.\]
由 5.6.1，$U^{H'}=U'$ \Term{smooth}{光滑}。考虑\Term{extension}{扩张}
\[1\longrightarrow U'\longrightarrow E'\longrightarrow H\longrightarrow1.\]
按构造，$H'$ 在 $E'$ 中\Term{center}{中心}，故\Term{normal-subgroup}{正规}。取商得\Term{exact-sequence}{正合列}
\[1\longrightarrow U'\longrightarrow E''\longrightarrow H''\longrightarrow1.\]
因 $U'$ \Term{smooth}{光滑}且 $P(H'')$ 成立，$H_1$、$H_2$ 在 $E''$ 中的像由某个 $u\in U'(k)$ \Term{conjugacy}{共轭}，于是 $H_1=\operatorname{int}(u)H_2$。
\end{proof}
为证明 5.1.1 ii) b)，注意因 $k$ \Term{algebraically-closed}{代数闭}，$H$ 有一条\Term{composition-series}{合成列}，其逐次商或者\Term{smooth}{光滑}，或者当 $p>0$ 时\Term{isomorphism}{同构}于 $\mu_p$。反复使用 5.7.2，便归结为 $H$ \Term{smooth}{光滑}或 $H=\mu_p$ 的情形。前者应用 5.1.1 ii) a)。余下 $H=\mu_p$。因 $U$ \Term{smooth}{光滑}，它有一条\Term{characteristic-composition-series}{特征合成列}，逐次商\Term{etale}{平展}或\Term{isomorphism}{同构}于 $(\mathbf G_a)^r$（3.9）。若 $U$ \Term{etale}{平展}，应用 5.5.1。因此最终只余 $H=\mu_p$、$U=\mathbf G_a^r$ 的情形。
须证明 $\mathrm H^1(\mu_p,U)=0$。这里不能应用 5.4.1 的方法，因为 $\mu_p$ 在 $(\mathbf G_a)^r$ 上的\Term{action}{作用}一般并非线性的。固定记号：$H'$ 为 $H$ 在 $E$ 中的一个\Term{lifting}{提升}，$\mathfrak e=\operatorname{Lie}E$、$\mathfrak u=\operatorname{Lie}U$、$\mathfrak h=\operatorname{Lie}H$、$\mathfrak h'=\operatorname{Lie}H'$。取非零 $X\in\mathfrak h$ 满足 $X^{(p)}=X$（附录 II 3.1），设 $X'$ 为它在 $\mathfrak h'$ 中的\Term{lifting}{提升}。
因 $\mu_p$ 为\Term{height}{高度} $1$ 的\Term{radicial}{根性}群，$H$ 在 $E$ 中的\Term{lifting}{提升}\Term{set}{集合}与如下\Term{set}{集合} $A$ 一一对应：$A$ 由\Term{projection}{投影}到 $X$ 且满足 $Y^{(p)}=Y$ 的 $Y\in\mathfrak e$ 构成（附录 II 2.2）。同样，若 $Y\in A$ 对应于\Term{lifting}{提升} $H''$，$u\in U(k)$，则 $\operatorname{int}(u)H'=H''$ 当且仅当 $\operatorname{Ad}(u)X=Y$。设 $B$ 为 $\mathfrak e$ 中形如 $\operatorname{Ad}(u)X$（$u\in U(k)$）的元素的\Term{set}{集合}。显然 $B\subset A$；只须证明 $k$ \Term{algebraically-closed}{代数闭}时 $A=B$。
\SourcePage{376}
\textbf{a) 研究 $A$。}因 $\mathfrak u$ 交换且被 $\mathfrak h$ \Term{normal-subgroup}{正规}化，Jacobson 公式（第 VII\textsubscript{A} 讲 5.2）在这里对 $u\in\mathfrak u$ 给出
\[(X'+u)^{(p)}=X'^{(p)}+u^{(p)}+(\operatorname{ad}X')^{p-1}(u)
=X'+(\operatorname{ad}X')^{p-1}(u),\]
因此 $X'+u\in A\Leftrightarrow u=(\operatorname{ad}X')^{p-1}(u)$。
设 $\mathfrak u=\coprod_{n\in\mathbf Z/p\mathbf Z}\mathfrak u_n$ 为 $\mathfrak u$ 在 $H'$ \Term{action}{作用}下的\Term{canonical}{典范}\Term{resolution}{分解}，也可写成
\[\mathfrak u=\mathfrak u_0\oplus\coprod_{n\in(\mathbf Z/p\mathbf Z)^\times}\mathfrak u_n.\]
若 $u\in\mathfrak u_0$，则 $\operatorname{ad}X'(u)=0$；若 $u\in\mathfrak u_n$，$n\in(\mathbf Z/p\mathbf Z)^\times$，则 $(\operatorname{ad}X')^{p-1}(u)=u$。最终，$Y=X'+u$ 属于 $A$ 当且仅当 $u\in\coprod_{n\in(\mathbf Z/p\mathbf Z)^\times}\mathfrak u_n$。因而 $A$ 是 $\mathbf W(\mathfrak e)$ 的一个\Term{irreducible}{不可约}\Term{subscheme}{子概形}的 $k$-\Term{rational-point}{有理点}集，其\Term{dimension}{维数}为
\[\operatorname{rg}\mathfrak u-\operatorname{rg}\mathfrak u_0
=\operatorname{rg}\mathfrak u-\operatorname{rg}\operatorname{Centr}_{\mathfrak u}(X').\]
\textbf{b) 研究 $B$。}需要以下\Term{lemma}{引理}。
\begin{SGAresult}{\Term{lemma}{引理}}{5.7.3（Rosenlicht）}
设 $U$ 为域 $k$ 上的\Term{unipotent}{幺幂}\Term{algebraic-group}{代数群}，\Term{action}{作用}于\Term{quasi-affine}{拟仿射} $k$-\Term{scheme}{概形} $X$。则每一点 $x\in X(k)$ 的\Term{orbit}{轨道}在 $X$ 中闭（这里的\Term{orbit}{轨道}指映射 $g\mapsto g\cdot x$ 下 $G$ 在 $\operatorname{ens}(X)$ 中的像）。
\end{SGAresult}
\begin{proof}
由\Term{functorial}{函子性}，$G$ \Term{action}{作用}于 $X$ 的\Term{affine-envelope}{仿射包络}，即 $\operatorname{Spec}\Gamma(X,\mathscr O_X)$，故可假定 $X$ \Term{affine}{仿射}。进而可假定 $k$ \Term{algebraically-closed}{代数闭}、$X$ \Term{reduced}{既约}、$U$ \Term{smooth}{光滑}（注意 $k$ \Term{perfect-field}{完美}时 $U_{\mathrm{red}}$ \Term{action}{作用}于 $X_{\mathrm{red}}$）。设 $Y$ 为 $G$ 在\Term{morphism}{态射} $g\mapsto g\cdot x$ 下的\Term{scheme-theoretic-image}{概形像}（EGA I 9.5.1），它是 $X$ 的闭\Term{reduced}{既约}\Term{subscheme}{子概形}，且 $G$ \Term{action}{作用}于其上。由 EGA IV 1.8.6 容易\Term{pushout}{推出}，$x$ 的\Term{orbit}{轨道}为 $Y$ 的\Term{dense}{稠密}开部分 $Z$。须证明 $Z=\operatorname{ens}(Y)$。设 $F$ 为以 $Y\setminus Z$ 为底\Term{space}{空间}的 $Y$ 的闭\Term{reduced}{既约}\Term{subscheme}{子概形}，则 $F=\mathrm V(J)$，其中 $J$ 为 $\Gamma(Y,\mathscr O_Y)$ 的非\Term{nilideal}{零理想}。因 $G$ \Term{smooth}{光滑}，它\Term{action}{作用}于 $F$，从而\Term{action}{作用}于 $J$，故 $J^G\ne0$（3.2）。若 $a\in J^G$ 非零，则 $a$ 在\Term{orbit}{轨道} $Z$ 上必为常数，因 $Z$ 在 $Y$ 中\Term{dense}{稠密}，它在 $Y$ 上亦为常数。于是\Term{ideal}{理想} $J$ 包含 $k$，故 $F=\varnothing$。
\end{proof}
将此\Term{lemma}{引理}用于通过\Term{adjoint-representation}{伴随表示}\Term{action}{作用}于\Term{affine}{仿射}\Term{space}{空间} $\mathbf W(\mathfrak e)$ 的群 $U$，得 $X'$ 的\Term{orbit}{轨道}是 $\mathbf W(\mathfrak e)$ 某个闭子\Term{prescheme}{预概形}的底\Term{space}{空间}。又 $X'$ 的\Term{stabilizer}{稳定子} $Z$ 是 $X'$ 在 $U$ 中的\Term{centralizer}{中心化子}，且有\Term{closed-immersion}{闭浸入}
\[U/Z\longrightarrow\mathbf W(\mathfrak e).\]
因此 $X'$ 的\Term{orbit}{轨道}为 $\mathbf W(\mathfrak e)$ 某个闭\Term{subscheme}{子概形}的底\Term{space}{空间}，其\Term{dimension}{维数}等于 $\dim U-\dim Z$。
\textbf{c) 完成 5.1.1 ii) b) 的证明。}当 $k$ \Term{algebraically-closed}{代数闭}时，\Term{canonical}{典范}映射 $U(k)\to(U/Z)(k)$ \Term{surjectivity}{满射}，故由 b)，$B$ 是 $\mathbf W(\mathfrak e)$ 某个闭\Term{subscheme}{子概形}的\Term{rational-point}{有理点}集，
\SourcePage{377}
\Term{dimension}{维数}为 $\dim U-\dim Z$。结合 a)，为证明 $A=B$，只须证明
\[\operatorname{rg}\mathfrak u-\operatorname{rg}\operatorname{Centr}_{\mathfrak u}(X')
\leq\dim U-\dim\operatorname{Centr}_U(X').\]
因 $U$ \Term{smooth}{光滑}，$\dim U=\operatorname{rg}\mathfrak u$；另一方面（第 II 讲 5.3.3），
\[\dim\operatorname{Centr}_U(X')\leq\operatorname{rg}\operatorname{Centr}_{\mathfrak u}(X'),\]
结论得证（事实上两者相等，这又证明了 $Z$ \Term{smooth}{光滑}（5.7.1））。

\textbf{5.8. 完成 5.1.1 i) 的证明。}还须证明 i) c) 和 i) d)。
\textbf{5.1.1 i) d) 的证明（$U$ \Term{smooth}{光滑}，$H$ \Term{connected}{连通}）。}需要以下\Term{lemma}{引理}。
\begin{SGAresult}{\Term{lemma}{引理}}{5.8.1}
沿用 5.1.1 的记号，假定 $U$ \Term{smooth}{光滑}、$H$ \Term{radicial}{根性}，且 $H_1$ 为 $H$ 的一个\Term{algebraic-subgroup}{代数子群}，在 $E$ 中有\Term{lifting}{提升} $H'_1$。则 $H$ 在 $E$ 中有包含 $H'_1$ 的\Term{lifting}{提升} $H'$。
\end{SGAresult}
\begin{proof}
对 $H/H_1$ 的\Term{height}{高度}归纳，可以假定 $H/H_1$ \Term{height}{高度}为 $1$。设 $C'=\operatorname{Centr}_E(H'_1)$。断言\Term{canonical}{典范}\Term{morphism}{态射} $C'\to H$ 为\Term{epimorphism}{满态射}。为此可以假定 $k$ \Term{algebraically-closed}{代数闭}；这时\Term{extension}{扩张} $E$ 平凡（5.1.1 i) a)）。设 $H''$ 为 $H$ 在 $E$ 中的一个\Term{lifting}{提升}，$H''_1$ 为 $H_1$ 在 $H''$ 中的\Term{inverse-image}{逆像}。群 $H'_1$ 和 $H''_1$ 为 $H_1$ 在 $E$ 中的两个\Term{lifting}{提升}，因 $k$ \Term{algebraically-closed}{代数闭}且 $U$ \Term{smooth}{光滑}，它们由 $U(k)$ 中某元素\Term{conjugacy}{共轭}（5.1.1 ii) b)）。显然只须为 $C''=\operatorname{Centr}_E(H''_1)$ 证明上述性质；但 $C''$ 包含 $H''$，故性质显然。又由 5.7.1，$C\cap U$ \Term{smooth}{光滑}。因此可以用 $C$ 代替 $E$，从而假定 $H'_1$ \Term{center}{中心}。再取模 $H'_1$ 的商，可假定 $H_1=0$、$H$ \Term{height}{高度}为 $1$。
因 $U$ \Term{smooth}{光滑}，有 $p$-\Term{lie-algebra}{Lie 代数}的\Term{exact-sequence}{正合列}（附录 II 3.2）
\[0\longrightarrow\operatorname{Lie}U\longrightarrow\operatorname{Lie}E\longrightarrow\operatorname{Lie}H\longrightarrow0.\tag{*}\]
由附录 II 2.2，$E$ 平凡\Term{equivalence}{等价}于 (*) 为平凡的 $p$-\Term{lie-algebra}{Lie 代数}\Term{extension}{扩张}。假定 $H\ne0$（故 $\operatorname{Lie}H\ne0$），并假定找到 $\mathfrak h=\operatorname{Lie}H$ 的一个非零 $p$-\Term{lie-subalgebra}{Lie 子代数} $\mathfrak h_1$，\Term{liftable}{可提升}为 $\operatorname{Lie}E$ 的 $p$-\Term{lie-subalgebra}{Lie 子代数} $\mathfrak h'_1$。由同处，有 $H$ 的\Term{subgroup}{子群} $H_1$ 使 $\operatorname{Lie}H_1=\mathfrak h_1$，且 $H_1$ 在 $E$ 中有\Term{lifting}{提升} $H'_1$，使 $\operatorname{Lie}H'_1=\mathfrak h'_1$。再次运用上述化简，便归结为将 $H$ 换为 $H/H_1$ 的同一问题。因 $H$ \Term{height}{高度}为 $1$，$\operatorname{Lie}(H/H_1)=\operatorname{Lie}H/\operatorname{Lie}H_1$（同处），故 $\operatorname{rg}\operatorname{Lie}(H/H_1)\leq\operatorname{rg}\operatorname{Lie}H-1$。总之，对 $\operatorname{Lie}H$ 的秩归纳可知，只须在 $\mathfrak h\ne0$ 时找到 $\operatorname{Lie}H$ 的一个\Term{liftable}{可提升}到 $\operatorname{Lie}E$ 中的非零 \Term{lie-subalgebra}{Lie 子代数}。为此有以下\Term{lemma}{引理}。
\end{proof}
\begin{SGAresult}{\Term{lemma}{引理}}{5.8.2}
设 $k$ 为\Term{characteristic}{特征} $p>0$ 的域，$\varphi:\mathfrak g\to\mathfrak h$ 为\Term{finite-rank}{有限秩} $p$-$k$-\Term{lie-algebra}{Lie 代数}的\Term{surjectivity}{满射}。则：
\SourcePage{378}
i) 若 $\mathfrak h_{\bar k}$ 为\Term{reductive-lie-algebra}{约化}的（4.2.2）且非零，则 $\mathfrak g$ 中有一个\Term{reductive-lie-algebra}{约化} \Term{lie-subalgebra}{Lie 子代数} $\mathfrak h_1$，其在 $\mathfrak h$ 中的像非零。
ii) 若 $k$ \Term{perfect-field}{完美}，$u\in\mathfrak h$ \Term{unipotent}{幺幂}（即存在 $n$ 使 $u^{(p^n)}=0$），则 $u$ \Term{liftable}{可提升}为 $\mathfrak g$ 的\Term{unipotent}{幺幂}元素。
\end{SGAresult}
\begin{proof}
在 i) 中取 $\mathfrak h$ 的非零元素 $X$，在 ii) 中取 $X=u$，设 $X'$ 为 $X$ 在 $\mathfrak g$ 中的\Term{lifting}{提升}。由 $X'$ 生成的 $\mathfrak g$ 的 $p$-\Term{lie-subalgebra}{Lie 子代数}是阿贝尔 $p$-\Term{lie-algebra}{Lie 代数} $\mathfrak h'$（第 VII 讲）。
情形 i)：由 4.2.2 的描述，$\mathfrak h'_{\bar k}$ 的\Term{reductive-lie-algebra}{约化}部分（同处）已定义在 $k$ 上，记为 $\mathfrak r'$。断言其在 $\mathfrak h$ 中的像非零。为此可假定 $k$ \Term{algebraically-closed}{代数闭}，故 $\mathfrak h'=\mathfrak r'\oplus\mathfrak u'$（$\mathfrak u'$ 为其\Term{unipotent}{幺幂}部分）。若 $\mathfrak r'$ 在 $\mathfrak h$ 中的像为零，则 $\mathfrak h'$ 在 $\mathfrak h$ 中的像\Term{unipotent}{幺幂}，因 $\mathfrak h$ \Term{reductive-lie-algebra}{约化}而必为零（参见 2.4 ii)），但按构造此像包含 $X$。
情形 ii)：交换 $\mathfrak r'$、$\mathfrak u'$ 的角色，同样处理。
\end{proof}
\textbf{完成 5.8.1 的证明。}设 $H\ne0$，由 5.8.2，$\operatorname{Lie}E$ 有非零\Term{reductive-lie-algebra}{约化} \Term{lie-subalgebra}{Lie 子代数} $\mathfrak h'_1$。因 $\operatorname{Lie}U$ \Term{unipotent}{幺幂}（4.3 i)），必有 $(\operatorname{Lie}U)\cap\mathfrak h'_1=0$，故 $\mathfrak h'_1$ 是 $\operatorname{Lie}H$ 某个 $p$-\Term{lie-subalgebra}{Lie 子代数}的\Term{lifting}{提升}。
\textbf{完成 5.1.1 i) d) 的证明。}由 5.8.1，有 $E$ 的\Term{algebraic-subgroup}{代数子群}族 $H'_n$（$n\in\mathbf N$），使 $H'_{n+1}\supset H'_n$，且 $H'_n$ \Term{lifting}{提升} ${}_{F^n}H$。递减\Term{subgroup}{子群}列 $\operatorname{Centr}_E(H'_n)$ \Term{stable}{稳定}，设其\Term{stable}{稳定}值为 $C$。$C$ 的\Term{center}{中心} $Z$ 包含所有 $H'_n$，故 $Z$ 在 $H$ 中的像包含所有 ${}_{F^n}H$，从而为 $H$ 的\Term{open-subgroup}{开子群}（第 VII\textsubscript{A} 讲 §4），因 $H$ \Term{connected}{连通}而等于 $H$。故为证明 $E$ 平凡，可以用 $Z$ 代替 $E$，从而假定 $E$ 交换。我们将在 7.2.1 中看到，$E$ 包含一个极大\Term{multiplicative-type}{乘法型}\Term{subgroup}{子群} $M$，其形成与\Term{base-field}{基域}\Term{extension}{扩张}交换。因 $E_{\bar k}$ 为平凡\Term{extension}{扩张}（5.1.1 i) a)）且 $U$ \Term{unipotent}{幺幂}，显然 $M$ 是 $H$ 在 $E$ 中的唯一\Term{lifting}{提升}。
\textbf{5.1.1 i) c) 的证明（$U$ 为 $k$-\Term{solvable}{可解}的）。}因 $(H/H')^0(k)$ 的阶与 $p$ 互素，由\Term{duality}{对偶性}立即可知，存在整数 $n$ 使 ${}_nH$ 为\Term{etale}{平展}\Term{subgroup}{子群}，且\Term{canonical}{典范}\Term{morphism}{态射} ${}_nH\to H/H^0$ 为\Term{epimorphism}{满态射}，从而 $H/{}_nH$ \Term{connected}{连通}。由 5.2.3 d)，${}_nH$ 在 $E$ 中有\Term{lifting}{提升} $H'$。与 5.8.1 证明开头一样，可证 $C=\operatorname{Centr}_E(H')$ 为 $E$ 的\Term{subgroup}{子群}，使 $C\to H$ 为\Term{epimorphism}{满态射}且 $C\cap U$ \Term{smooth}{光滑}。以 $C$ 代替 $E$，再取模 $H'$ 的商，便归结为 $U$ \Term{smooth}{光滑}、$H$ \Term{connected}{连通}的情形，即 i) d)。这完成了 5.1.1 的证明。

\textbf{5.9. 反例。}先给出构造以\Term{unipotent}{幺幂} $k$-群 $U$ 为核的\Term{multiplicative-type}{乘法型} $k$-群 $H$ 的非平凡\Term{extension}{扩张}的一种方法。给定 $H$ 在 $U$ 上的\Term{action}{作用}，设 $E=U\cdot H$ 为\Term{semidirect-product}{半直积}。再取由 $1$-\Term{cocycle}{余循环} $a$ 表示的 $\mathrm H^1(k,U)$ 中一个元素。群 $U$
\SourcePage{379}
以\Term{inner-automorphism}{内自同构}\Term{action}{作用}于 $E$，故 $a$ 定义 $E$ 的一个 $k$-\Term{form}{形式}，记为 $E_a$。再假定 $U$ 交换，则 $U$ 在 $U$ 及商 $E/U=H$ 上的\Term{action}{作用}平凡，故 $E_a$ 仍是以 $U$ 为核的 $H$ 的\Term{extension}{扩张}。为简便，假定 $H$ 为\Term{etale}{平展}\Term{diagonalizable}{可对角化}群，因而 $\bar k/k$ 的 Galois 群 $\mathscr G$ 在 $H(\bar k)$ 上\Term{trivial-action}{平凡作用}；$\mathscr G$ 在 $E_a(\bar k)$ 的点上的\Term{action}{作用}由下式给出：
\[{}^g(u,h)=({}^gu+a(g)-{}^ha(g),h)
\qquad(g\in\mathscr G,\ u\in U(\bar k),\ h\in H(\bar k)).\]
若 $h\in H(k)$，$X$ 为其在 $E_a$ 中的\Term{inverse-image}{逆像}，则 $X$ 为由 $\mathscr G$ 取值于 $U$ 的 $1$-\Term{cocycle}{余循环} $g\mapsto a(g)-{}^ha(g)$ 的类定义的 $U$-\Term{torsor}{挠子}。因此，只要存在 $h\in H(k)$ 使此 $1$-\Term{cocycle}{余循环}非平凡，\Term{extension}{扩张} $E_a$ 就非平凡。将此构造用于以下两个特例。
\textbf{a) 以 $\mathbf Z/p\mathbf Z$ 为核的\Term{etale}{平展}\Term{diagonalizable-group}{可对角化群} $H$ 的非\Term{trivial-extension}{平凡扩张}。}
取 $H=(\mathbf Z/p\mathbf Z)^\times$，以乘法\Term{action}{作用}于 $U=\mathbf Z/p\mathbf Z$。再设 $k$ 为\Term{characteristic}{特征} $p$ 的域，$K/k$ 为 Galois 群 $\mathscr G\simeq\mathbf Z/p\mathbf Z$ 的\Term{extension}{扩张}，$a$ 为 $\mathrm H^1(\mathscr G,U)=\operatorname{Hom}(\mathscr G,U)$ 的非零元素。群 $E_a$ 即为所求。
\textbf{b) 以\Term{smooth}{光滑}\Term{connected}{连通}\Term{unipotent-group}{幺幂群} $U$ 为核的\Term{etale}{平展}\Term{diagonalizable-group}{可对角化群} $H$ 的非\Term{trivial-extension}{平凡扩张}。}
取非\Term{perfect-field}{完美}域 $k$，使存在 $\mathbf G_a$ 的一个 $k$-\Term{form}{形式} $U$，满足 $\mathrm H^1(k,U)\ne0$。例如（参见 J.-P. Serre，《Cohomologie Galoisienne》），可取 $k$ 为\Term{equal-characteristic}{等特征} $p>0$ 的\Term{discrete-valuation-ring}{离散赋值环}的\Term{fraction-field}{分式域}，$U$ 为 $\mathbf G_a\times_k\mathbf G_a$ 中由方程 $X^p+X+tY^p=0$ 定义的\Term{algebraic-subgroup}{代数子群}，其中 $t$ 是 $k$ 的一个\Term{uniformizer}{一致化参数}。为简便，再假定 $k$ 包含所有 $(p-1)$ 次\Term{root-of-unity}{单位根}，则有\Term{algebraic-group}{代数群}\Term{exact-sequence}{正合列}
\[0\longrightarrow\mathbf Z/p\mathbf Z\longrightarrow U\longrightarrow\mathbf G_a\longrightarrow0,
\qquad(x,y)\longmapsto y,\]
从而有\Term{cohomology}{上同调}\Term{exact-sequence}{正合列}
\[\mathbf G_a(k)\xrightarrow d\mathrm H^1(k,\mathbf Z/p\mathbf Z)\longrightarrow\mathrm H^1(k,U)\longrightarrow0,\]
其中 $d$ 将 $x$ 对应于方程 $X^p+X+tx^p=0$ 定义的 $\mathbf Z/p\mathbf Z$-\Term{homogeneous-space}{齐性空间}。又已知 $\mathrm H^1(k,\mathbf Z/p\mathbf Z)\simeq k/\mathfrak P(k)$，其中 $\mathfrak P(x)=x^p+x$，故
\[\mathrm H^1(k,U)\simeq k/(\mathfrak P(k)+tk^p).\]
再假定 $p\ne2$，则显然 $t^2\notin\mathfrak P(k)+tk^p$，故 $\mathrm H^1(k,U)\ne0$。另一方面，$\mu_{p-1}$ 按 $(h,x,y)\mapsto(hx,hy)$ \Term{action}{作用}于 $U$。以 $\mathscr G$ 表示方程 $X^p+X+t^2=0$ 所定义的\Term{extension}{扩张} $K$ 的 Galois 群，
\SourcePage{380}
设 $a\in\mathrm H^1(\mathscr G,U)$ 为上述非零元素。直接验证可知，此时 $E_a$ 是以 $U$ 为核的 $\mu_{p-1}$ 的非平凡\Term{extension}{扩张}。
\textbf{c) 以 $\alpha_p$ 为核的 $\mathbf G_m$ 的非\Term{trivial-extension}{平凡扩张}。}
由 5.1.1 i) b)，这种\Term{extension}{扩张}只可能存在于非\Term{perfect-field}{完美}域 $k$ 上。设 $k$ 非\Term{perfect-field}{完美}，$G$ 为 $U=\mathbf G_a$ 与 $H=\mathbf G_m$ 的\Term{semidirect-product}{半直积}，其中 $\mathbf G_m$ 以\Term{homothety}{位似}\Term{action}{作用}于 $\mathbf G_a$。因 $U$ 在 $G$ 中\Term{normal-subgroup}{正规}，${}_FU$ 也\Term{normal-subgroup}{正规}。令 $G'=G/{}_FU$。群 $G'$ \Term{isomorphism}{同构}于 $\mathbf G_a\cdot\mathbf G_m$，其中 $\mathbf G_m$ 按 $(h,u)\mapsto h^pu$ \Term{action}{作用}于 $\mathbf G_a$。$G$（相应地，$G''$）的子\Term{torus}{环面}\Term{functor}{函子} $\mathscr T_G$（相应地，$\mathscr T_{G''}$）（参见第 XV 讲）\Term{isomorphism}{同构}于 $\mathbf G_a$，且由 $G\to G''$ \Term{derive}{导出}的\Term{morphism}{态射} $\mathscr T_G\to\mathscr T_{G''}$ 等同于 $u\mapsto u^p$。因此，若 $T''$ 为 $G''$ 中对应于点 $x\in k\simeq\mathbf G_a(k)$ 的子\Term{torus}{环面}，且 $x^{1/p}\notin k$，则 $T''$ 在 $G$ 中的\Term{inverse-image}{逆像} $E$ 是以 ${}_FU=\alpha_p$ 为核的\Term{torus}{环面} $T''$ 的\Term{extension}{扩张}，没有定义在 $k$ 上的极大\Term{torus}{环面}，故并非平凡。计算得 $E$ 是 $G=\operatorname{Spec}[U,T,T^{-1}]$ 中由方程 $U^p=x-xT^p$ 定义的\Term{subgroup}{子群}。
\textbf{注记 5.9.1。}此例说明，定义在非\Term{perfect-field}{完美}域上的非\Term{smooth}{光滑}\Term{algebraic-group}{代数群}未必有极大\Term{torus}{环面}，从而回答了第 XIV 讲 1.5 b) 所提的问题。
\textbf{d) 下面给出以\Term{unipotent-group}{幺幂群} $U$ 为核的\Term{multiplicative-type}{乘法型群} $H$ 的\Term{trivial-extension}{平凡扩张} $E$，其中有 $H$ 的两个\Term{lifting}{提升} $H'$、$H''$ 不由 $U(k)$ 中任何元素\Term{conjugacy}{共轭}。}
取 $E$ 为 $U=\mathbf G_a$ 与 $\mu_p=\operatorname{Spec}k[T]/(T^p-1)$ 的\Term{semidirect-product}{半直积}，$\mu_p$ 在 $\mathbf G_a=\operatorname{Spec}k[U]$ 上的\Term{action}{作用}由\Term{comorphism}{余态射}
\[U\longmapsto(U+U^p)T-U^p\]
定义。$\mu_p$ 在 $U$ 中的\Term{centralizer}{中心化子}是方程 $U+U^p=0$ 定义的\Term{etale}{平展}群 $Z$。因此，若 $k$ 不\Term{algebraically-closed}{代数闭}，\Term{canonical}{典范}映射 $U(k)\to(U/Z)(k)$ 一般不\Term{surjectivity}{满射}，故由 5.1.1 ii) b)，$\mu_p$ 在 $E$ 中的两个\Term{lifting}{提升}未必由 $U(k)$ 中某个元素\Term{conjugacy}{共轭}。
另一个例子如下，这次 $k$ 是\Term{characteristic}{特征} $p>0$ 的\Term{algebraically-closed}{代数闭}域。设 $G$ 为 $\mu_p$ 与 $\alpha_p$ 的\Term{semidirect-product}{半直积}\Term{radicial}{根性}群，其中 $\mu_p$ 以“\Term{homothety}{位似}”\Term{action}{作用}于 $\alpha_p$。有带 $\mu_p$ \Term{operator}{算子}群的\Term{algebraic-group}{代数群}\Term{exact-sequence}{正合列}
\[0\longrightarrow\alpha_p\longrightarrow\mathbf G_a\xrightarrow{x\mapsto x^p}\mathbf G_a\longrightarrow0,\]
其中 $\mu_p$ 在第一个 $\mathbf G_a$ 上以\Term{homothety}{位似}\Term{action}{作用}，在第二个上\Term{trivial-action}{平凡作用}。\Term{cohomology}{上同调}\Term{exact-sequence}{正合列}（附录 I，命题 11）在此给出
\[0\longrightarrow\mathbf G_a(k)\longrightarrow\mathrm H^1(\mu_p,\alpha_p)\longrightarrow\mathrm H^1(\mu_p,\mathbf G_a).\]
最后一项为零（I 5.3.3），故 $\mathrm H^1(\mu_p,\alpha_p)$ 非零，因而 $\mu_p$ 在 $G$ 中的两个\Term{lifting}{提升}未必\Term{conjugacy}{共轭}。

\SourcePage{381}
\SGAsection{6}{以乘法型群为核的幺幂群扩张}
\textbf{6.1. \Term{theorem}{定理}的陈述。}
\begin{SGAresult}{\Term{theorem}{定理}}{6.1.1}
设 $k$ 为域，$U$ 为\Term{unipotent}{幺幂} $k$-\Term{algebraic-group}{代数群}，$H$ 为\Term{multiplicative-type}{乘法型} $k$-群，$E$ 为以 $H$ 为核的 $U$ 的 $k$-\Term{algebraic-group}{代数群}\Term{extension}{扩张}，即有\Term{exact-sequence}{正合列}
\[1\longrightarrow H\longrightarrow E\longrightarrow U\longrightarrow1.\]
则在下列各情形中，\Term{extension}{扩张} $E$ 平凡，且 $U$ 在 $E$ 中有唯一\Term{lifting}{提升}：
\begin{enumerate}
\item[A)] $U$ \Term{smooth}{光滑}，且下列条件之一成立：
\begin{enumerate}
\item[i)] $U$ \Term{connected}{连通}，且\Term{canonical}{典范}\Term{morphism}{态射} $E\to U$ 有\Term{section}{截面}；
\item[ii)] $U$ 有一条逐次商\Term{isomorphism}{同构}于 $\mathbf G_a$ 的\Term{composition-series}{合成列}；
\item[iii)] $H$ \Term{etale}{平展}；
\item[iv)] $k$ \Term{perfect-field}{完美}。
\end{enumerate}
\item[B)] $U=\alpha_p$，且 $k$ \Term{perfect-field}{完美}。
\item[C)] $E$ 交换，且 $k$ \Term{perfect-field}{完美}。
\end{enumerate}
\end{SGAresult}

\textbf{6.2. 6.1.1 A) 的证明。}先证明三个\Term{lemma}{引理}。
\begin{SGAresult}{\Term{lemma}{引理}}{6.2.1}
设 $S$ 为\Term{prescheme}{预概形}，$E$ 为 $S$-\Term{group-prescheme}{群预概形}，它是以\Term{multiplicative-type}{乘法型}且\Term{finite-type}{有限型}的 $S$-群 $H$ 为核的 $S$-\Term{group-prescheme}{群预概形} $U$ 的\Term{extension}{扩张}，且 $U$ 的各\Term{fiber}{纤维}\Term{connected}{连通}（即 $U$ 是 $E$ 对 $H$ 的 fpqc 商）。则 $E$ 是\Term{central-extension}{中心扩张}。
\end{SGAresult}
\begin{proof}
因 $H$ 交换，群 $U$ 通过一个 $S$-群\Term{morphism}{态射}
\[u:U\longrightarrow\operatorname{Aut}_{S\text{-gr}}(H)\]
以\Term{inner-automorphism}{内自同构}\Term{action}{作用}于 $H$。\Term{functor}{函子} $\operatorname{Aut}_{S\text{-gr}}(H)$ 由\Term{etale}{平展} $S$-\Term{scheme}{概形}表示（第 X 讲 5.10），因而其\Term{unit-section}{单位截面}为既开又闭的\Term{immersion}{浸入}。因 $U$ 的各\Term{fiber}{纤维}\Term{connected}{连通}，\Term{pushout}{推出} $u$ 是单位\Term{morphism}{态射}。
\end{proof}
\begin{SGAresult}{\Term{lemma}{引理}}{6.2.2}
沿用 6.1.1 的记号。若 $E$ 平凡，则 $U$ 在 $E$ 中有唯一\Term{lifting}{提升}。
\end{SGAresult}
\begin{proof}
设 $U'$、$U''$ 为 $U$ 在 $E$ 中的两个\Term{lifting}{提升}。为证明 $U'=U''$，可以假定 $k$ \Term{algebraically-closed}{代数闭}，只须证明 $\mathrm H^1(U,H)=0$。若 $U$ \Term{connected}{连通}，则 $U$ \Term{centralizes}{中心化} $H$（6.2.1），故由 2.4 ii)，
\[\mathrm H^1(U,H)=\operatorname{Hom}_{k\text{-gr}}(U,H)=0.\]
一般情形中，设 $U'_1$ 为 $U$ 的\Term{connected-component}{连通分支} $U^0$ 在 $E$ 中的唯一\Term{lifting}{提升}，令 $N=\operatorname{Norm}_E(U'_1)$。若 $g\in E(k)$，则 $\operatorname{int}(g)U'_1$ 是 $U^0$ 的\Term{lifting}{提升}，故等于 $U'_1$，从而 $N(k)\supset E(k)$。又 $N$ 包含 $H$（6.2.1）及 $U'_1$，立即得 $N=E$。取模 $U'_1$ 的商，便归结为 $U$ \Term{etale}{平展}的情形。此时因 $k$ \Term{algebraically-closed}{代数闭}，$\mathrm H^1(U,H)$
\SourcePage{382}
等同于普通\Term{cohomology-group}{上同调群} $\mathrm H^1(U(k),H(k))$\footnote{编者注：例如参见 M. Demazure 和 P. Gabriel，《Groupes algébriques I》，Masson \& North-Holland (1970) 的命题 III.6.4.2。}，故为零，因为 $U(k)$ 为\Term{finite}{有限} $p$-群，$H(k)$ 则唯一 $p$-可除。
\end{proof}
\begin{SGAresult}{推论}{6.2.3}
为证明 $E$ 平凡，只须证明它在\Term{base-field}{基域}的一个\Term{finite}{有限}\Term{separable}{可分}\Term{extension}{扩张}后平凡；特别地，可以假定 $H$ \Term{diagonalizable}{可对角化}。
\end{SGAresult}
\begin{SGAresult}{\Term{lemma}{引理}}{6.2.4}
若以 $\mathbf G_m$ 为核的 $U$ 的每个\Term{central-extension}{中心扩张}都平凡，则以任意\Term{diagonalizable}{可对角化} $k$-群 $H$ 为核的 $U$ 的每个\Term{central-extension}{中心扩张} $E$ 都平凡。
\end{SGAresult}
\begin{proof}
先对 $r$ 归纳，知所作假设蕴涵 $H=\mathbf G_m^r$ 时 $E$ 平凡。一般情形，$H$ 可\Term{embedding}{嵌入}某个 $\mathbf G_m^r$（由\Term{duality}{对偶性}立即可知）。设 $H''=\mathbf G_m^r/H$，得到\Term{exact-sequence}{正合列}（附录 I 2.1）
\[\mathrm Z^1(U,H'')\longrightarrow\operatorname{Ext}_{\mathrm{alg}}(U,H)\longrightarrow\operatorname{Ext}_{\mathrm{alg}}(U,\mathbf G_m^r)=0,\]
其中 $U$ 在 $H$、$H''$、$\mathbf G_m^r$ 上\Term{trivial-action}{平凡作用}。但 $\mathrm Z^1(U,H'')=\operatorname{Hom}_{k\text{-gr}}(U,H'')=0$（2.4 ii)），故 $\operatorname{Ext}_{\mathrm{alg}}(U,H)=0$。
\end{proof}
\textbf{6.1.1 A) i) 的证明。}因 $E\to U$ 有\Term{section}{截面}，$E$ 定义 $\mathrm H^2(U,H)$ 中的元素 $f$（附录 I 3.1）。须证明 $f$ 为零；只须证明每个 $2$-\Term{cocycle}{余循环} $f:U\times_kU\to H$ 都是常值\Term{morphism}{态射}。这将由下述\Term{lemma}{引理}得出。
\begin{SGAresult}{\Term{lemma}{引理}}{6.2.5}
设 $U$ 为\Term{smooth}{光滑}\Term{connected}{连通}\Term{unipotent}{幺幂} $k$-\Term{algebraic-group}{代数群}，$H$ 为\Term{multiplicative-type}{乘法型} $k$-群。则任意 $k$-\Term{morphism}{态射}（\Term{prescheme}{预概形}\Term{morphism}{态射}）$f:U\to H$ 都是常值的。
\end{SGAresult}
\begin{proof}
可以假定 $k$ \Term{algebraically-closed}{代数闭}，对 $\dim U$ 递增归纳。若 $\dim U>0$，则 $U$ 有\Term{composition-series}{合成列}（参见 3.9）
\[1\longrightarrow U'\longrightarrow U\xrightarrow\pi U''\longrightarrow1,\]
其中 $U'\simeq\mathbf G_a$，$\dim U''<\dim U$。只须证明 $f$ 通过 $U''$ \Term{resolution}{分解}。因 $\pi$ 定义的\Term{equivalence-relation}{等价关系}的图在 $k$ 上\Term{smooth}{光滑}，故\Term{reduced}{既约}，只须证明：若 $x,y\in U(k)$ 在 $U''(k)$ 中有相同像 $z$，则 $f(x)=f(y)$。但 $\pi^{-1}(z)\simeq\mathbf G_a$，故 $f$ 在其上的\Term{restriction}{限制}通过 $H$ 的一个\Term{reduced}{既约}\Term{irreducible-component}{不可约分支}\Term{resolution}{分解}，从而通过一个\Term{isomorphism}{同构}于 $(\mathbf G_m)^r$ 的 $k$-\Term{scheme}{概形}\Term{resolution}{分解}。最后只须注意，从 $\mathbf G_a$ 到 $(\mathbf G_m)^r$ 的每个\Term{morphism}{态射}均常值，因为 $\mathbf G_a$ 上每个\Term{invertible}{可逆}\Term{regular-function}{正则函数}都为常数。
\end{proof}
\textbf{6.1.1 A) ii) 的证明。}由 6.2.1、6.2.3 和 6.2.4，可以假定 $H=\mathbf G_m$。按假设，$U$ 有\Term{composition-series}{合成列} $U\supset U_1\supset\cdots$，各 $U_i/U_{i+1}$ \Term{isomorphism}{同构}于 $\mathbf G_a$。设 $E_1$ 为 $U_1$ 在 $E$ 中的\Term{inverse-image}{逆像}。对 $\dim U$ 归纳，可以假定\Term{extension}{扩张} $E_1$ 平凡；设 $U'_1$ 为 $U_1$ 的唯一
\SourcePage{383}
\Term{lifting}{提升}。与 6.2.4 的证明一样，$U'_1$ 在 $E$ 中\Term{normal-subgroup}{正规}。取模 $U'_1$ 的商后，便归结为 $U=\mathbf G_a$ 的情形。此时 $S$-\Term{scheme}{概形} $E$（$S=U$）为 $S$-群 $\mathbf G_m\times_kS$ 下的\Term{torsor}{挠子}，因 $\operatorname{Pic}(S)=0$，它有\Term{section}{截面}（第 VIII 讲 4.3）。于是由 A) i)，\Term{extension}{扩张} $E$ 平凡。
\textbf{6.1.1 A) iii) 的证明（$U$ \Term{smooth}{光滑}，$H$ \Term{etale}{平展}）。}先设 $U$ \Term{connected}{连通}。群 $U_{\bar k}$ 有一条逐次商\Term{isomorphism}{同构}于 $\mathbf G_a$ 的\Term{composition-series}{合成列}，故由 A) ii)，$E_{\bar k}$ 平凡。因 $H$ \Term{etale}{平展}，$U_{\bar k}$ 在 $E_{\bar k}$ 中的唯一\Term{lifting}{提升}显然为 $E_{\bar k}$ 的\Term{connected-component}{连通分支}，故该\Term{lifting}{提升}已定义在 $k$ 上。一般情形中，取模 $E$ 的\Term{connected-component}{连通分支}的商，便归结为 $E$ \Term{etale}{平展}的情形，再由 6.2.3 归结为 $E$ \Term{completely-split}{完全分裂}的情形。因 $U(k)$ 为 $p$-群，$H(k)$ 的阶与 $p$ 互素，可以取 $E(k)$ 的 Sylow $p$-\Term{subgroup}{子群}作为 $U(k)$ 的\Term{lifting}{提升}。
\textbf{6.1.1 A) iv) 的证明（$U$ \Term{smooth}{光滑}，$k$ \Term{perfect}{完美}）。}若 $U$ \Term{connected}{连通}，则它有一条逐次商\Term{isomorphism}{同构}于 $\mathbf G_a$ 的\Term{composition-series}{合成列}（4.1.2 b)），应用 A) ii) 即可。一般情形中，由上面的结果可归结为 $U$ \Term{etale}{平展}，再归结为 $U$ \Term{completely-split}{完全分裂}、$H$ \Term{diagonalizable}{可对角化}的情形（6.2.3）。利用 $H$ 的\Term{characteristic-composition-series}{特征合成列} $H\supset H^0\supset{}_{F^n}H\supset0$，便归结为 $H$ 属于下列三类之一：a) $H$ \Term{etale}{平展}；b) $H=\mathbf G_m^r$；c) $H$ \Term{radicial}{根性}。
a) 应用 A) iii)；c) 应用 1.6。最后在 b) 中，由 Hilbert \Term{theorem}{定理} 90，$E\to U$ 有\Term{section}{截面}，故只须证明 $\mathrm H^2(U,\mathbf G_m^r)=\mathrm H^2(U(k),\mathbf G_m^r(k))=0$。$U(k)$ 为\Term{finite}{有限} $p$-群，而 $\mathbf G_m^r(k)$ 唯一 $p$-可除（因 $k$ \Term{perfect-field}{完美}），故所求成立。

\textbf{6.3. 6.1.1 B) 和 C) 的证明。}由 6.2.1、6.2.3、6.2.4，只须在 $H=\mathbf G_m$ 时证明 B)。有\Term{exact-sequence}{正合列}
\[1\longrightarrow\mathbf G_m\longrightarrow E\longrightarrow\alpha_p\longrightarrow1.\]
因 $\mathbf G_m$ \Term{smooth}{光滑}，得到 $p$-\Term{lie-algebra}{Lie 代数}的\Term{exact-sequence}{正合列}（附录 II 3.2）
\[0\longrightarrow\operatorname{Lie}\mathbf G_m\longrightarrow\operatorname{Lie}E\longrightarrow\operatorname{Lie}\alpha_p\longrightarrow0.\tag{*}\]
因 $\alpha_p$ \Term{height}{高度}为 $1$，\Term{extension}{扩张} $E$ 平凡当且仅当 $p$-\Term{lie-algebra}{Lie 代数}\Term{exact-sequence}{正合列} (*) \Term{splitting}{分裂}（附录 II 2.2）。已知 $\operatorname{Lie}\alpha_p$ 由满足 $X^{(p)}=0$ 的非零元素 $X$ 生成（附录 II 2.1），故只须证明 $X$ \Term{liftable}{可提升}为 $\operatorname{Lie}E$ 中满足 $Z^{(p)}=0$ 的元素 $Z$。由 5.8.2 ii)，有 $\operatorname{Lie}E$ 的\Term{unipotent}{幺幂}元素 $Z$ \Term{lifting}{提升} $X$。因 $\operatorname{Lie}E$ 的\Term{unipotent}{幺幂}部分显然\Term{dimension}{维数}至多为 $1$，必有 $Z^{(p)}=0$。
\textbf{6.1.1 C) 的证明。}若 $U'_1$ 为 $U$ 的\Term{algebraic-subgroup}{代数子群} $U_1$ 在 $E$ 中的\Term{lifting}{提升}，因假设 $E$ 交换，$U'_1$ 在 $E$ 中\Term{normal-subgroup}{正规}，故可取模 $U'_1$ 的商。利用 $U$ 的\Term{composition-series}{合成列}（3.5 ii)），
\SourcePage{384}
上述讨论便将问题归结为 $U$ \Term{smooth}{光滑}或 $U=\alpha_p$ 的情形。此时由 A) iv) 和 B)，$E$ 平凡。

\textbf{6.4. 以\Term{multiplicative-type}{乘法型群} $H$ 为核的\Term{unipotent-group}{幺幂群} $U$ 的非\Term{trivial-extension}{平凡扩张}的例子。}
由 6.1.1 A) iv)，此问题只会在\Term{characteristic}{特征} $p>0$ 时出现。
\textbf{a) $H=\mathbf G_m$，$U$ 为 $\mathbf G_a$ 的非平凡\Term{form}{形式}（参见附录 III §5）。}
设 $k$ 为\Term{characteristic}{特征} $2$ 的非\Term{perfect-field}{完美}域，$u\in k$ 满足 $u^{1/2}\notin k$。考虑环为 $k[X,Y,(X^2+uY^2)^{-1}]$ 的\Term{affine}{仿射}群 $E$，其乘法由\Term{comorphism}{余态射}
\[(X,Y)\longmapsto(XX'+uYY',XY'+YX')\]
给出。群 $E$ \Term{smooth}{光滑}、\Term{connected}{连通}、交换，\Term{dimension}{维数}为 $2$；\Term{subscheme}{子概形} $Y=0$ 定义\Term{subgroup}{子群} $H\simeq\mathbf G_m$。$E$ 中平方映射的核 $K$ 由方程 $X^2+uY^2=1$ 给出，故\Term{dimension}{维数}为 $1$。群 $K$ 包含 $E$ 的\Term{unipotent-radical}{幺幂根基}（定义在 $\bar k$ 上），也包含来自 $H$ 的\Term{isomorphism}{同构}于 $\mu_2$ 的部分。因 $K$ 在 $k$ 上\Term{reduced}{既约}，$E$ 的\Term{unipotent-radical}{幺幂根基}未定义在 $k$ 上，$U=E/H$ 不能\Term{lifting}{提升}到 $E$。（直接验证可知，$U$ 为 $\mathbf G_a$ 的一个\Term{form}{形式}，其环为 $k[V,W]/(V+uV^2+W^2)$，\Term{morphism}{态射} $E\to U$ 对应于\Term{comorphism}{余态射}
\[V\longmapsto\frac{Y^2}{X^2+Y^2u},\qquad W\longmapsto\frac{XY}{X^2+Y^2u}.\]）
\textbf{b) $H=\mathbf G_m$，$U=\mathbf Z/2\mathbf Z$，$k$ 为\Term{characteristic}{特征} $2$ 的非\Term{perfect-field}{完美域}。}
按 a) 选取 $k$、$u$，考虑 $\operatorname{GL}_2$ 中由下述元素 $X$ 生成的\Term{subgroup}{子群} $E$：
\[X=\begin{pmatrix}0&1\\u&0\end{pmatrix}
=\begin{pmatrix}u^{1/2}&0\\0&u^{1/2}\end{pmatrix}
\begin{pmatrix}0&u^{-1/2}\\u^{1/2}&0\end{pmatrix}.\]
群 $E$ 为以 $\mathbf G_m$ 为核的 $\mathbf Z/2\mathbf Z$ 的\Term{extension}{扩张}；但因 $X$ 的\Term{unipotent}{幺幂}部分未定义在 $k$ 上，此\Term{extension}{扩张}并不平凡。
\textbf{c) $H=\mu_p$，$U=\alpha_p$，$k$ 非\Term{perfect}{完美}。}
设 $\mathfrak e$ 为由 $X,Y$ 生成的交换 $p$-\Term{lie-algebra}{Lie 代数}，满足 $X^{(p)}=X$、$Y^{(p)}=aX$。由附录 II 2.2，$\mathfrak e$ 是以 $\mu_p$ 为核的 $\alpha_p$ 的\Term{algebraic-group}{代数群}\Term{extension}{扩张} $E$ 的 $p$-\Term{lie-algebra}{Lie 代数}。但此\Term{extension}{扩张}平凡当且仅当存在 $b\in k$ 使 $b^p=a$（这时 $(bX+Y)^{(p)}=0$）。
\textbf{d) $H=\mu_2$，$U=\alpha_2\times\alpha_2$，$E$ 不交换，$k$ 的\Term{characteristic}{特征}为 $2$。}
考虑特殊\Term{linear-group}{线性群} $\operatorname{SL}_{2,k}$，令 $E={}_F(\operatorname{SL}_{2,k})$。群 $E$ 为\Term{height}{高度} $1$ 的\Term{radicial}{根性}群，其 \Term{lie-algebra}{Lie 代数}由 $X,Y,Z$ 生成，满足
\[[X,Y]=Z,\qquad[X,Z]=[Y,Z]=0,\]
\[X^{(p)}=Y^{(p)}=0,\qquad Z^{(p)}=Z.\]
\SourcePage{385}
因此 $E$ 为以 $\mu_2$ 为核的 $U\simeq\alpha_2\times\alpha_2$ 的\Term{central-extension}{中心扩张}。$U$ 的每个因子 $\alpha_2$ 都唯一地\Term{lifting}{提升}到 $E$ 中；但因 $[X,Y]\ne0$，$U$ 本身不能\Term{lifting}{提升}到 $E$。

\SGAsection{7}{幂零仿射代数群}
\textbf{7.1. \Term{multiplicative-type}{乘法型群}的\Term{group-extension}{扩张}。}
\begin{SGAresult}{命题}{7.1.1}
设 $S$ 为\Term{prescheme}{预概形}，$H$、$K$ 为两个\Term{multiplicative-type}{乘法型}且\Term{finite-type}{有限型}的 $S$-\Term{group-prescheme}{群预概形}，$E$ 为以 $H$ 为核的 $K$ 的\Term{group-prescheme}{群预概形}\Term{extension}{扩张}（即 $K$ 为 $E$ 对 $H$ 的 fpqc 商）。则在下列两种情形中，$E$ 为\Term{multiplicative-type}{乘法型}：a) $E$ 交换；b) $K$ 的各\Term{fiber}{纤维}\Term{connected}{连通}。
\end{SGAresult}
\begin{proof}
i) $S$ 为域 $k$ 的谱的情形。所证断言为 fpqc \Term{local-scheme}{局部}的，故可以假定 $k$ \Term{algebraically-closed}{代数闭}，从而 $H$、$K$ \Term{diagonalizable}{可对角化}。注意 $K$ 以\Term{inner-automorphism}{内自同构}在 $H$ 上\Term{trivial-action}{平凡作用}：a) 中显然，b) 中来自 6.2.1。对 $K$ 的适当\Term{composition-series}{合成列}的\Term{length}{长度}归纳，便归结为 $K$ 属于下列三类之一：a) $K=\mathbf G_m$；b) $K=\mu_p$；c) $K=\mu_q$、$(q,p)=1$，且这时 $E$ 交换。再利用 $H$ 到 $\mathbf G_m^r$ 的\Term{embedding}{嵌入}，便得 $E$ \Term{embedding}{嵌入}以 $\mathbf G_m^r$ 为核的 $K$ 的一个\Term{extension}{扩张}，因而可假定 $H=\mathbf G_m^r$。
a) 若 $K=\mathbf G_m$，则 $E$ 是\Term{smooth}{光滑}、\Term{connected}{连通}、\Term{affine}{仿射}（第 VI\textsubscript{B} 讲 9.2）的\Term{algebraic-group}{代数群}，其\Term{unipotent}{幺幂}秩为零，故为\Term{torus}{环面}。
b) $K=\mu_p$。此时 $E$ 为平凡\Term{extension}{扩张}。事实上，因 $H$ \Term{smooth}{光滑}，\Term{canonical}{典范}\Term{morphism}{态射} $\operatorname{Lie}E\to\operatorname{Lie}\mu_p$ \Term{surjectivity}{满射}（附录 II 3.2），结合附录 II 2.2，应用 5.8.2 i) 即可。
c) $K=\mu_q$、$(q,p)=1$，且 $E$ 交换。此时\Term{extension}{扩张} $E$ 也平凡。事实上，设 $x\in E(k)$ \Term{lifting}{提升} $\mu_q(k)$ 的一个\Term{generator}{生成元} $\bar x$。元素 $x^q\in H(k)$，故形如 $y^q$，$y\in H(k)$（注意 $\mathbf G_m^r(k)$ 为 $q$-可除的）。因 $E$ 交换，$y^{-1}x$ 是 $\bar x$ 的一个阶为 $q$ 的\Term{lifting}{提升}。
ii) 一般情形。$H$、$K$ 在 $S$ 上\Term{flat}{平坦}、\Term{affine}{仿射}且\Term{finitely-presented}{有限表现}（第 IX 讲 2.1），故 $E$ 也如此（第 VI\textsubscript{B} 讲 9.2）。运用第 VI\textsubscript{B} 讲 §10 的一般技术，归结为 $S$ \Term{noetherian}{诺特}。为证明 $E$ 为\Term{multiplicative-type}{乘法型}，只须证明 $E$ 交换，且每个 ${}_nE$ 都在 $S$ 上\Term{finite}{有限}（第 X 讲 4.8 b)）。
a) $E$ 交换。须验证\Term{morphism}{态射}
\[E\times_SE\longrightarrow E,\qquad(x,y)\longmapsto[x,y]=xyx^{-1}y^{-1}\]
通过 $E$ 的\Term{unit-section}{单位截面}\Term{resolution}{分解}；只须在 $S$ 为\Term{artinian-local-ring}{局部阿廷环}的谱时验证。但此时由 i) 和第 X 讲 2.3，$E$ 为\Term{multiplicative-type}{乘法型}。
\SourcePage{386}
b) ${}_nE$ 在 $S$ 上\Term{finite}{有限}。事实上，有\Term{exact-sequence}{正合列}
\[0\longrightarrow{}_nH\longrightarrow{}_nE\longrightarrow{}_nK\xrightarrow u H_n,\]
其中 $H_n$ 为 $H$ 中 $n$ 次幂映射的\Term{cokernel}{余核}。已知 $H_n$ 为\Term{multiplicative-type}{乘法型}（第 IX 讲 2.7），故\Term{separated}{分离}，且 ${}_nH$、${}_nK$ 在 $S$ 上\Term{finite}{有限}。$\operatorname{Ker}u$ 为 ${}_nK$ 的\Term{closed-subgroup}{闭子群}，故在 $S$ 上\Term{finite}{有限}；$\!{}_nE$ 为\Term{finite-group}{有限群}对\Term{finite-group}{有限群}的\Term{extension}{扩张}，故在 $S$ 上\Term{finite}{有限}（第 VI\textsubscript{B} 讲 9.2）。
\end{proof}

\textbf{7.2. 交换\Term{affine-algebraic-group}{仿射代数群}的结构。}
\begin{SGAresult}{\Term{theorem}{定理}}{7.2.1}
设 $k$ 为域，$G$ 为交换\Term{affine}{仿射} $k$-\Term{algebraic-group}{代数群}。则：
a) $G$ 包含最大\Term{multiplicative-type}{乘法型}\Term{subgroup}{子群} $M$。$M$ 是 $G$ 的\Term{characteristic-subgroup}{特征子群}，$G/M$ \Term{unipotent}{幺幂}，且 $M$ 的形成与域 $k$ 的\Term{extension}{扩张}交换。
b) 若 $k$ \Term{perfect-field}{完美}，则 $G$ 唯一地为 $M$ 与一个\Term{unipotent}{幺幂}\Term{algebraic-subgroup}{代数子群} $U$ 的\Term{direct-product}{直积}。
\end{SGAresult}
\begin{proof}
i) $k$ \Term{algebraically-closed}{代数闭}。若 $G$ \Term{smooth}{光滑}，7.2.1 b) 是熟知的（BIBLE §4，\Term{theorem}{定理} 4）。若 $G$ 为\Term{height}{高度} $1$ 的\Term{radicial}{根性}群，由 4.3.1 v) 和附录 II 2.2，4.2.2 中 $\operatorname{Lie}G$ 的\Term{resolution}{分解}对应于 $G$ 的一个如 7.2.1 b) 的\Term{resolution}{分解}。一般情形中，$G$ 有一条逐次商\Term{smooth}{光滑}或\Term{radicial}{根性}且\Term{height}{高度}为 $1$ 的\Term{composition-series}{合成列}（附录 II 3.1）。为证明 7.2.1 b)，只须注意：若有交换\Term{algebraic-group}{代数群}\Term{exact-sequence}{正合列}
\[0\longrightarrow G'\longrightarrow G\longrightarrow G''\longrightarrow0,\]
其中 $G'$（相应地，$G''$）是\Term{multiplicative-type}{乘法型}群与\Term{unipotent}{幺幂}群的乘积 $G'=M'\cdot U'$（相应地，$G''=M''\cdot U''$），则 $G$ 也如此。事实上，考虑\Term{exact-sequence}{正合列}
\[0\longrightarrow G'/U'\longrightarrow G/U'\longrightarrow G''\longrightarrow0.\]
由 7.1.1 a)，$M''$ 在 $G/U'$ 中的\Term{inverse-image}{逆像}是\Term{multiplicative-type}{乘法型}\Term{subgroup}{子群} $M_1$。群 $M_1$ \Term{liftable}{可提升}为 $G$ 中的\Term{subgroup}{子群} $M$（5.1.1 i) a)）。同样，这次用 6.1.1 C)，可证 $G$ 有\Term{unipotent}{幺幂}\Term{subgroup}{子群} $U$，为以 $U'$ 为核的 $U''$ 的\Term{extension}{扩张}。显然 $G=M\cdot U$。
ii) $k$ 任意。由 i)，$G_{\bar k}$ 为\Term{multiplicative-type}{乘法型}群 $M_{\bar k}$ 与\Term{unipotent}{幺幂}群 $U_{\bar k}$ 的\Term{direct-product}{直积}。令 $S=\bar k\times_k\bar k$，设 $M_1$、$M_2$ 为 $M_{\bar k}$ 沿两\Term{projection}{投影} $S\rightrightarrows\bar k$ 的\Term{inverse-image}{逆像}。群 $G_S/M_2=(G/M_{\bar k})_S$ 的各\Term{fiber}{纤维}\Term{unipotent}{幺幂}，故 $M_1$ 在其中的像为零（2.4 i)），从而 $M_2\supset M_1$。同样 $M_1\supset M_2$，最终 $M_1=M_2$。由 fpqc \Term{descent}{下降}，$M_{\bar k}$ 来自 $G$ 的一个\Term{algebraic-subgroup}{代数子群} $M$。显然 $M$ 为\Term{multiplicative-type}{乘法型}、$G/M$ \Term{unipotent}{幺幂}，且 $M$ 的形成与 $k$ 的任意\Term{extension}{扩张}相容。对每个 $k$-\Term{prescheme}{预概形} $S$，$G_S$ 的每个\Term{multiplicative-type}{乘法型}\Term{subgroup}{子群} $H$ 都包含于 $M_S$。事实上，由 2.5，它在各\Term{fiber}{纤维}\Term{unipotent}{幺幂}的群 $(G/M)_S=G_S/M_S$ 中的像为零。特别地，
\SourcePage{387}
取 $H$ 为 $M_S$ 在 $G_S$ 的一个\Term{automorphism}{自同构}下的变换，就\Term{pushout}{推出} $M$ 为 $G$ 的\Term{characteristic-subgroup}{特征子群}。
最后，若 $k$ \Term{perfect-field}{完美}，由 6.1.1 C)，$G/M$ 可唯一地\Term{lifting}{提升}为 $G$ 中的\Term{unipotent}{幺幂}群。
\end{proof}
\textbf{注记 7.2.2。}
i) 若 $k$ 非\Term{perfect-field}{完美}，$G_{\bar k}$ 的\Term{unipotent}{幺幂}分支未必定义在 $k$ 上，例 6.4 a) 即说明这一点。
ii) 与\Term{multiplicative-type}{乘法型}分支 $M$ 的情形不同，\Term{unipotent}{幺幂}分支 $U$ 一般并非 $G$ 的\Term{characteristic-subgroup}{特征子群}（无论 $k$ 的\Term{characteristic}{特征}如何）。当然，\Term{resolution}{分解} 7.2.1 b) 的唯一性保证 $U$ 在 $G$ 的每个 $k$-\Term{automorphism}{自同构}下\Term{invariant}{不变}。但若 $U$、$M$ 使得存在一个 $k$-\Term{prescheme}{预概形} $S$ 及非零 $S$-\Term{homomorphism}{同态} $h:U_S\to M_S$（参见 2.6），则得到 $G_S$ 的 $S$-\Term{automorphism}{自同构} $(u,m)\mapsto(u,h(u)+m)$，它并不保持 $U_S$ \Term{invariant}{不变}。
iii) 若 $G$ 在 $k$ 上\Term{finite}{有限}，则在 Cartier \Term{duality}{对偶性}（2.6）下，$G/M$ 对应于 $G$ 的\Term{duality}{对偶} $D(G)$ 的\Term{connected-component}{连通分支}。

\textbf{7.3. \Term{nilpotent}{幂零}\Term{affine-algebraic-group}{仿射代数群}的结构。}
\begin{SGAresult}{\Term{theorem}{定理}}{7.3.1}
设 $k$ 为域，$G$ 为\Term{nilpotent}{幂零}（第 VI\textsubscript{B} 讲 §8）、\Term{affine}{仿射}、\Term{connected}{连通}的 $k$-\Term{algebraic-group}{代数群}。则 $G$ 有最大\Term{multiplicative-type}{乘法型}\Term{subgroup}{子群} $M$。群 $M$ \Term{center}{中心}且为\Term{characteristic-subgroup}{特征子群}，$G/M$ 为\Term{unipotent}{幺幂}\Term{algebraic-group}{代数群}。
\end{SGAresult}
设 $Z$ 为 $G$ 的\Term{center}{中心}，$M$ 为 $Z$ 的最大\Term{multiplicative-type}{乘法型}\Term{subgroup}{子群}（7.2.1）。因 $Z$ 为 $G$ 的\Term{characteristic-subgroup}{特征子群}，$M$ 为 $Z$ 的\Term{characteristic-subgroup}{特征子群}，故 $M$ 为 $G$ 的\Term{characteristic-subgroup}{特征子群}。只须证明 $G/M$ \Term{unipotent}{幺幂}。对 $G$ 的\Term{ascending-central-series}{升中心列}\Term{length}{长度}归纳，此结论来自以下更一般的\Term{lemma}{引理}。
\begin{SGAresult}{\Term{lemma}{引理}}{7.3.2（Rosenlicht）}
设 $G$ 为\Term{connected}{连通} $k$-\Term{algebraic-group}{代数群}，$Z$ 为其\Term{center}{中心}。则 $G'=G/Z$ 的\Term{center}{中心} $Z'$ \Term{unipotent}{幺幂}。
\end{SGAresult}
\begin{proof}
可假定 $k$ \Term{algebraically-closed}{代数闭}。只须证明：对每个素数 $\ell$，$Z'$ 都不含\Term{isomorphism}{同构}于 $\mu_\ell$ 的\Term{subgroup}{子群}（4.6.1 vi)）。设 $\mu_\ell$ 为 $Z'$ 的\Term{subgroup}{子群}，$N$ 为其在 $G$ 中的\Term{inverse-image}{逆像}。因 $\mu_\ell$ 在 $G'$ 中\Term{center}{中心}，$N$ 在 $G$ 中\Term{normal-subgroup}{正规}。
i) $(\ell,p)=1$ 的情形。可以找到 $x\in N(k)$ 和整数 $n$，使：a) $x$ \Term{lifting}{提升} $\mu_\ell$ 的一个\Term{generator}{生成元} $\bar x$；b) $x^{\ell^n}\in Z^0(k)$。（只须选择 $\bar x$ 的一个\Term{lifting}{提升}，使它在 $N/Z^0(k)$ 中的像属于 Sylow $\ell$-\Term{subgroup}{子群}。）
\Term{abelian-group}{交换群} $Z^0$ 中的 $\ell$ 次幂映射为\Term{etale}{平展}\Term{morphism}{态射}，故 $Z^0(k)$ 为 $\ell$-可除的。因此将 $x$ 乘以 $Z^0(k)$ 的某个元素后，可以假定 $x^{\ell^n}=0$。群 $N$ 由两个彼此交换的\Term{abelian-group}{交换群}
\SourcePage{388}
（$Z$ 及由 $x$ 生成的群）生成，故交换。群 ${}_{\ell^n}N$ 为\Term{multiplicative-type}{乘法型}，是 $N$ 的\Term{characteristic-subgroup}{特征子群}，故在 $G$ 中\Term{normal-subgroup}{正规}；因 $G$ \Term{connected}{连通}，它\Term{center}{中心}（第 IX 讲 5.5）。于是 ${}_{\ell^n}N$ 包含于 $Z$，这与它在 $G'$ 中的像包含 $\mu_\ell$ 矛盾。
ii) $\ell=p$。存在整数 $n$，使 ${}_{F^n}N$ 在 $G'$ 中的像包含 $\mu_p$，于是有\Term{exact-sequence}{正合列}
\[1\longrightarrow K\longrightarrow{}_{F^n}N\longrightarrow\mu_p\longrightarrow1.\tag{*}\]
群 $K$ 包含于 $Z$，故交换；由 7.2.1、7.1.1 b) 和 5.5.1，存在 ${}_{F^n}N$ 中的一个\Term{multiplicative-type}{乘法型}\Term{subgroup}{子群}，其在模 $K$ 的商中的像为 $\mu_p$。与 i) 一样\Term{pushout}{推出} ${}_{F^n}N$ 交换。${}_{F^n}N$ 的\Term{multiplicative-type}{乘法型}分支（7.2.1）为 ${}_{F^n}N$ 的\Term{characteristic-subgroup}{特征子群}，故在 $G$ 中\Term{normal-subgroup}{正规}，因而\Term{center}{中心}（第 IX 讲 5.5）；但它在 $G'$ 中的像包含 $\mu_p$，又得矛盾。
\end{proof}

\SGAsection{A}{附录 I：Hochschild 上同调与代数群扩张}
\textbf{A.1. \Term{cohomology-group}{上同调群}的定义。}设 $k$ 为域，$G$ 为 $k$-\Term{algebraic-group}{代数群}，$(\mathrm{Ab})_G$ 为带 $G$ \Term{action}{作用}的交换 $k$-\Term{algebraic-group}{代数群}所成的\Term{abelian-category}{阿贝尔范畴}。若 $A\in\operatorname{Ob}(\mathrm{Ab})_G$，则 $A$ \Term{canonical}{典范}定义的\Term{functor}{函子} $h_A:(\mathrm{Sch}/k)^\circ\to(\mathrm{Ens})$ 为 I 3.2 意义下的 $G$-$\mathbf Z$-模。因此可以考虑 $G$ 取值于 $A$ 的\Term{algebra}{代数}\Term{cochain}{余链}的\Term{standard-complex}{标准复形} $\mathrm C^\bullet(G,A)$（第 I 讲 5.1），以及 $i$-\Term{cocycle}{余循环}群 $\mathrm Z^i(G,A)$、$i$-\Term{coboundary}{余边界}群 $\mathrm B^i(G,A)$ 和第 $i$ 个\Term{cohomology-group}{上同调群} $\mathrm H^i(G,A)$。如常，$\mathrm H^0(G,A)$ 等同于 $A^G(k)$，其中 $A^G$ 为 $A$ 在 $G$ 下的\Term{invariants}{不变量} $k$-\Term{functor}{函子}。群 $\mathrm H^1(G,A)$ 分类带 $G$ \Term{action}{作用}的平凡 $A$-\Term{principal-homogeneous-space}{主齐性空间}。
\Term{functor}{函子} $A\mapsto\mathrm H^\bullet(G,A)$ 一般并非从\Term{category}{范畴} $(\mathrm{Ab})_G$ 到 $\mathrm{Ab}$ 的\Term{cohomology}{上同调}\Term{functor}{函子}；但有下述命题。
\begin{SGAresult}{命题}{A.1.1}
设 $0\to A'\xrightarrow u A\xrightarrow v A''\to0$ 为 $(\mathrm{Ab})_G$ 中的\Term{exact-sequence}{正合列}。则：
a) 若 $v$ 有\Term{section}{截面}（即有 $k$-\Term{prescheme}{预概形}\Term{morphism}{态射} $s:A''\to A$ 满足 $vs=1_{A''}$），则有通常的\Term{cohomology}{上同调}\Term{exact-sequence}{正合列}
\[\cdots\longrightarrow\mathrm H^i(G,A)\longrightarrow\mathrm H^i(G,A'')\xrightarrow d\mathrm H^{i+1}(G,A')\longrightarrow\cdots.\]
b) 若 $A(k)\to A''(k)$ \Term{surjectivity}{满射}，则有\Term{exact-sequence}{正合列}
\[0\to A'^G(k)\to A^G(k)\to A''^G(k)\xrightarrow d\mathrm H^1(G,A')\to\mathrm H^1(G,A)\to\mathrm H^1(G,A'').\tag{1}\]
\end{SGAresult}
\begin{proof}
a) \Term{section}{截面}的存在保证如下\Term{complex}{复形}列\Term{exact}{正合}：
\[0\longrightarrow\mathrm C^\bullet(G,A')\longrightarrow\mathrm C^\bullet(G,A)\longrightarrow\mathrm C^\bullet(G,A'')\longrightarrow0.\]
\SourcePage{389}
b) 若 $x''\in A''^G(k)$，它在 $A$ 中的\Term{inverse-image}{逆像}是一个平凡的 $A'$-\Term{principal-homogeneous-space}{主齐性空间}（因假定 $A(k)\to A''(k)$ \Term{surjectivity}{满射}），且带有 $G$ \Term{action}{作用}，故定义元素 $d(x'')\in\mathrm H^1(G,A')$。列 (1) 的\Term{exactness}{正合性}随即得到。
\end{proof}

\textbf{A.2. 群 $\operatorname{Ext}_{\mathrm{alg}}(G,A)$。}设 $A$、$G$ 为两个 $k$-\Term{algebraic-group}{代数群}，$E$、$E'$ 为以 $A$ 为核的 $G$ 的两个 $k$-\Term{algebraic-group}{代数群}\Term{extension}{扩张}。若有 $k$-群\Term{morphism}{态射} $u:E\to E'$ 使下图交换，就称这两个\Term{extension}{扩张}\Term{isomorphism}{同构}：
\[
\begin{tikzcd}
 & E\arrow[dr]\arrow[dd,"u"] & \\
 A\arrow[ur]\arrow[dr] & &G\\
 &E'\arrow[ur]&
\end{tikzcd}
\]
群 $E$ 以\Term{inner-automorphism}{内自同构}\Term{action}{作用}于 $A$；若 $A$ 交换，此\Term{action}{作用}通过 $G$ \Term{resolution}{分解}，从而 $G$ \Term{action}{作用}于 $A$。反之，若 $A\in\operatorname{Ob}(\mathrm{Ab})_G$，则以 $\operatorname{Ext}_{\mathrm{alg}}(G,A)$ 表示以 $A$ 为核的 $G$ 的\Term{algebra}{代数}\Term{extension}{扩张} $E$ 的类所成的\Term{set}{集合}，其中要求由 $E$ 定义的 $G$ 在 $A$ 上的\Term{action}{作用}，与 $A$ 的 $(\mathrm{Ab})_G$ \Term{object}{对象}结构所给的\Term{action}{作用}一致。
$\operatorname{Ext}_{\mathrm{alg}}(G,A)$ 自然地为关于 $A$ \Term{covariant}{协变}、关于 $G$ \Term{contravariant}{反变}的\Term{bifunctor}{双函子}。具体地：
a) 若 $f:A\to B$ 为 $(\mathrm{Ab})_G$ 中的\Term{morphism}{态射}，$E$ 表示 $\operatorname{Ext}_{\mathrm{alg}}(G,A)$ 的一个元素，则定义 $f_*(E)\in\operatorname{Ext}_{\mathrm{alg}}(G,B)$ 为以下\Term{extension}{扩张}的类：将\Term{semidirect-product}{半直积} $B\cdot E$（$E$ 通过 $G$ \Term{action}{作用}于 $B$）对 $A$ 在\Term{morphism}{态射} $(f,i)$ 下的像\Term{algebraic-subgroup}{代数子群}取商（由第 VI\textsubscript{A} 讲 §5，此商\Term{representable}{可表示}）。于是有\Term{commutative-diagram}{交换图}
\[
\begin{tikzcd}
1\arrow[r]&A\arrow[r,"i"]\arrow[d,"f"]&E\arrow[r]\arrow[d]&G\arrow[r]\arrow[d,"1_G"]&1\\
1\arrow[r]&B\arrow[r]&f_*(E)\arrow[r]&G\arrow[r]&1.
\end{tikzcd}
\]
b) 若 $g:H\to G$ 为 $k$-\Term{algebraic-group}{代数群}的 $k$-\Term{morphism}{态射}，$E$ 为以 $A$ 为核的 $G$ 的\Term{extension}{扩张}，则\Term{fiber-product}{纤维积} $E\times_GH$ 自然为以 $A$ 为核的 $H$ 的\Term{extension}{扩张}，记为 $g^*(E)$。于是有\Term{commutative-diagram}{交换图}
\[
\begin{tikzcd}
1\arrow[r]&A\arrow[r]\arrow[d,"1_A"]&g^*(E)\arrow[r]\arrow[d]&H\arrow[r]\arrow[d,"g"]&1\\
1\arrow[r]&A\arrow[r]&E\arrow[r]&G\arrow[r]&1.
\end{tikzcd}
\]
适当调整 J.-P. Serre，《Groupes algébriques et corps de classe》第 VII 章的证明，可赋予 $\operatorname{Ext}_{\mathrm{alg}}(G,A)$ 自然的\Term{abelian-group}{阿贝尔群}结构，它关于 $A$、$G$ 都具有\Term{functorial}{函子性}。
\SourcePage{390}
\begin{SGAresult}{命题}{A.2.1}
设 $0\to A'\to A\to A''\to0$ 为 $(\mathrm{Ab})_G$ 中的\Term{exact-sequence}{正合列}，则：
a) 有\Term{abelian-group}{阿贝尔群}的\Term{canonical}{典范}\Term{exact-sequence}{正合列}
\[\mathrm Z^1(G,A)\to\mathrm Z^1(G,A'')\xrightarrow d\operatorname{Ext}_{\mathrm{alg}}(G,A')\to\operatorname{Ext}_{\mathrm{alg}}(G,A)\to\operatorname{Ext}_{\mathrm{alg}}(G,A'').\]
b) 若 $A(k)\to A''(k)$ \Term{surjectivity}{满射}，由 a) 得\Term{exact-sequence}{正合列}
\[\mathrm H^1(G,A)\to\mathrm H^1(G,A'')\xrightarrow d\operatorname{Ext}_{\mathrm{alg}}(G,A')\to\operatorname{Ext}_{\mathrm{alg}}(G,A)\to\operatorname{Ext}_{\mathrm{alg}}(G,A'').\tag{2}\]
\end{SGAresult}
\Term{exact-sequence}{正合列} a) 推广了交换\Term{extension}{扩张}情形中通常的 $\operatorname{Hom}(\cdot,\cdot)$ \Term{exact-sequence}{正合列}（同处），其证明也相同。这里只回顾\Term{coboundary}{余边界} $d:\mathrm Z^1(G,A'')\to\operatorname{Ext}_{\mathrm{alg}}(G,A')$ 的定义。考虑由\Term{exact-sequence}{正合列} $0\to A'\to A\to A''\to0$ 显然\Term{derive}{导出}的\Term{extension}{扩张}
\[1\longrightarrow A'\longrightarrow A\cdot G\longrightarrow A''\cdot G\longrightarrow1.\]
若 $u\in\mathrm Z^1(G,A'')$，则按通常方法，$u$ 定义一个\Term{homomorphism}{同态}\Term{section}{截面} $u:G\to A''\cdot G$。于是 $d(u)=u^*(A\cdot G)$。

\textbf{A.3. $\mathrm H^2(G,A)$ 与 $\operatorname{Ext}_{\mathrm{alg}}(G,A)$ 的比较。}在抽象群的情形，熟知\Term{abelian-group}{阿贝尔群} $\mathrm H^2(G,A)$ 与 $\operatorname{Ext}(G,A)$ 之间存在\Term{functorial}{函子性}的\Term{isomorphism}{同构}。在这里同样，若 $A\in(\mathrm{Ab})_G$，则每个 $2$-\Term{cocycle}{余循环} $u\in\mathrm Z^2(G,A)$ 都对应于\Term{prescheme}{预概形} $A\times_kG$ 上的一个\Term{algebraic-group}{代数群}结构，使之成为 $\operatorname{Ext}_{\mathrm{alg}}(G,A)$ 中的元素。此\Term{extension}{扩张}平凡当且仅当 $u\in\mathrm B^2(G,A)$（参见第 III 讲 1.2.2）。回顾 $A\times G$ 上的\Term{composition-law}{合成律}由下式定义：
\[(a,g)(a',g')=a+{}^ga'+u(g,g').\]
显然，如此得到的以 $A$ 为核的 $G$ 的\Term{extension}{扩张}并非任意\Term{extension}{扩张}，因为它们有\Term{section}{截面}。反之，若 $E\in\operatorname{Ext}_{\mathrm{alg}}(G,A)$ 有\Term{section}{截面} $s$，则 $E$ \Term{isomorphism}{同构}于与以下 $2$-\Term{cocycle}{余循环} $u$ 关联的以 $A$ 为核的 $G$ 的\Term{extension}{扩张}：
\[u(g,g')=s(g)s(g')s(gg')^{-1}.\]
最终得到下述命题。
\begin{SGAresult}{命题}{A.3.1}
下列取值于\Term{abelian-group}{阿贝尔群}的\Term{bifunctor}{双函子}之间有\Term{functorial}{函子性}的\Term{isomorphism}{同构}：
\[(G,A)\longmapsto\mathrm H^2(G,A),\qquad(G,A)\longmapsto\operatorname{Ext}_s(G,A),\]
其中 $\operatorname{Ext}_s(G,A)$ 表示 $\operatorname{Ext}_{\mathrm{alg}}(G,A)$ 中由有\Term{section}{截面}的、以 $A$ 为核的 $G$ 的\Term{extension}{扩张}类构成的\Term{subgroup}{子群}。
\end{SGAresult}

\SGAsection{B}{附录 II：根性群的回顾与补充}
设 $p>1$ 为素数，$S$ 为 $\mathbf F_p$-\Term{prescheme}{预概形}。
\SourcePage{391}
\textbf{B.1. Frobenius \Term{morphism}{态射}。}对每个 $S$-\Term{prescheme}{预概形} $X$ 及整数 $n>0$，以 $P^n:X\to X$ 表示与 $\mathscr O_X$ 中的 $p^n$ 次幂映射对应的 $\mathbf F_p$-\Term{endomorphism}{自同态}，以 $X^{(n)}$ 表示 $X$ 沿\Term{morphism}{态射} $P^n:S\to S$ 的\Term{inverse-image}{逆像} $S$-\Term{prescheme}{预概形}。则存在唯一 $S$-\Term{morphism}{态射} $F^n:X\to X^{(n)}$ 使下图交换：
\[
\begin{tikzcd}[row sep=large,column sep=large]
&&X\arrow[dll,"P^n"']\arrow[dl,"F^n"]\arrow[ddl,bend left=15]\\
X\arrow[d]&X^{(n)}\arrow[l]\arrow[d]&\\
S&S\arrow[l,"P^n"]&
\end{tikzcd}
\]
显然 $F^n$ 等同于 $F^1=F$ 的“第 $n$ 次迭代”；$F$ 称为 $X/S$ 的 Frobenius \Term{morphism}{态射}。
若 $G$ 为 $S$-\Term{group-prescheme}{群预概形}，则 $G^{(n)}$ 也是 $S$-\Term{group-prescheme}{群预概形}，$F^n:G\to G^{(n)}$ 为 $S$-群\Term{morphism}{态射}。其核 ${}_{F^n}G$ 为 $G$ 的\Term{characteristic-subgroup}{特征子群}\Term{prescheme}{预概形}（即在\Term{functor}{函子} $\operatorname{Aut}_{S\text{-gr}}(G)$ 下\Term{stable}{稳定}），且在 $S$ 上\Term{radicial}{根性}。若 $G$ 为\Term{radicial}{根性} $S$-\Term{group-prescheme}{群预概形}，当 ${}_{F^h}G=G$ 时，称 $G$ 的\Term{height}{高度} $\leq h$。

\textbf{B.2. 群与 $p$-\Term{lie-algebra}{Lie 代数}。}若 $G$ 为 $S$-\Term{group-prescheme}{群预概形}，则 $\operatorname{Lie}(G)$（第 II 讲）自然为\Term{restriction}{限制} $p$-\Term{lie-algebra}{Lie 代数}（第 VII\textsubscript{A} 讲 §§5、6）。特别地有如下结果（参见第 VII\textsubscript{A} 讲）。
\begin{SGAresult}{命题}{B.2.1}
i) $\operatorname{Lie}(\alpha_p)_S=\operatorname{Lie}(\mathbf G_a)_S$\footnote{编者注：此处删去了 $\alpha_p$ 的定义，因为本讲开头已经给出。}为 $\mathscr O_S$-\Term{sheaf-of-modules}{模层}，在 $\mathscr O_S$ 上自由，秩为 $1$，由满足 $X^{(p)}=0$ 的元素 $X$ 生成。
ii) $\operatorname{Lie}(\mu_p)_S=\operatorname{Lie}(\mathbf G_m)_S$ 为 $\mathscr O_S$-\Term{sheaf-of-modules}{模层}，在 $\mathscr O_S$ 上自由，秩为 $1$，有\Term{canonical-basis}{典范基} $X$ 满足 $X^{(p)}=X$。
\end{SGAresult}
回顾第 VII\textsubscript{A} 讲 §7 证明的基本结果。
\begin{SGAresult}{\Term{theorem}{定理}}{B.2.2}
假定 $S$ \Term{affine}{仿射}，环为 $A$。则\Term{functor}{函子} $G\mapsto\operatorname{Lie}G$ 给出以下两个\Term{category}{范畴}之间的\Term{equivalence-of-categories}{范畴等价}：一边为在 $S$ 上\Term{finitely-presented}{有限表现}、\Term{flat}{平坦}、\Term{height}{高度}为 $1$，且 \Term{lie-algebra}{Lie 代数}为\Term{locally-free}{局部自由} $\mathscr O_S$-模的 $S$-\Term{group-prescheme}{群预概形} $G$ 所成的\Term{category}{范畴}；
\SourcePage{392}
另一边为\Term{locally-free}{局部自由}\Term{finite-rank}{有限秩}\Term{restriction}{限制} $p$-$A$-\Term{lie-algebra}{Lie 代数}所成的\Term{category}{范畴}。
此外，若 $G$ 如上，$H$ 为\Term{finitely-presented}{有限表现}的 $S$-\Term{group-prescheme}{群预概形}，则\Term{canonical}{典范}\Term{morphism}{态射}
\[\operatorname{Hom}_{S\text{-gr}}(G,H)\longrightarrow
\operatorname{Hom}_{p\text{-}A\text{-Lie}}(\operatorname{Lie}G(S),\operatorname{Lie}H(S))\]
为\Term{isomorphism}{同构}。
\end{SGAresult}

\textbf{B.3. \Term{radicial}{根性}群与\Term{smooth}{光滑}群。}现在假定 $S$ 为\Term{characteristic}{特征} $p$ 的域 $k$ 的谱。
回顾第 VI\textsubscript{A} 讲 §5 已证明：若 $G$ 为 $k$-\Term{algebraic-group}{代数群}，$H$ 为 $G$ 的\Term{algebraic-subgroup}{代数子群}，则\Term{sheaf}{层} $G/H$（fpqc \Term{topology}{拓扑}下的\Term{sheaf}{层}）\Term{representable}{可表示}。再回顾（VII\textsubscript{A} 8.3）：
\begin{SGAresult}{命题}{B.3.1}
设 $G$ 为 $k$-\Term{algebraic-group}{代数群}。则存在整数 $m$，使对所有 $n\geq m$，\Term{algebraic-group}{代数群} $G/{}_{F^n}G$ 都在 $k$ 上\Term{smooth}{光滑}。
\end{SGAresult}
\begin{SGAresult}{命题}{B.3.2}
考虑 $k$-\Term{algebraic-group}{代数群}\Term{exact-sequence}{正合列} $1\to G'\to G\xrightarrow uG''\to1$，以及下列断言：
i) \Term{morphism}{态射} $u$ \Term{smooth}{光滑}；ii) $G'$ 在 $k$ 上\Term{smooth}{光滑}；iii) 对每个整数 $n>0$，有\Term{exact-sequence}{正合列}
\[1\longrightarrow{}_{F^n}G'\longrightarrow{}_{F^n}G\longrightarrow{}_{F^n}G''\longrightarrow1;\]
iv) \Term{morphism}{态射} ${}_FG\to{}_FG''$ 为\Term{epimorphism}{满态射}；v) \Term{morphism}{态射} $(\operatorname{Lie}G)(k)\to(\operatorname{Lie}G'')(k)$ \Term{surjectivity}{满射}。
则有蕴涵关系
\[\text{i)}\Longleftrightarrow\text{ii)}\Longrightarrow\text{iii)}\Longrightarrow\text{iv)}\Longleftrightarrow\text{v)}.\]
此外，若 $G$ \Term{smooth}{光滑}，五条断言\Term{equivalence}{等价}。
\end{SGAresult}
\begin{proof}
i)$\Leftrightarrow$ii) 来自第 VI\textsubscript{B} 讲 9.2 vii)。
ii)$\Rightarrow$iii)：若 $G'$ \Term{smooth}{光滑}，则 $F^n:G'\to G'^{(n)}$ 为\Term{epimorphism}{满态射}，iii) 由“蛇形”图得到。
iii)$\Rightarrow$iv) 显然。iv)$\Leftrightarrow$v) 来自\Term{theorem}{定理} B.2.2。
$G$ \Term{smooth}{光滑}时 v)$\Rightarrow$ii)：此时 $G''$ \Term{smooth}{光滑}，v) 蕴涵 $\dim G'=\dim_k(\operatorname{Lie}G')(k)$，故 $G'$ \Term{smooth}{光滑}。
\end{proof}

\SGAsection{C}{附录 III：关于第 XV、XVI、XVII 讲的注记与补充}
\textbf{C.1。}第 XV 讲的命题 1.2 和 1.2 bis 或许在去掉 $H_0$ 在 $S_0$ 上\Term{smooth}{光滑}的假设后仍成立。特别地，当 $G$ \Term{finite}{有限}、\Term{flat}{平坦}、交换时确实如此。
\SourcePage{393}
\textbf{C.2. 对 XV 4.8 的补充。}下述命题及后面的\Term{theorem}{定理} C.3.1、C.4.1 将发表于 M. Raynaud 正在撰写的关于\Term{discrete-valuation-ring}{离散赋值环}上的\Term{group-scheme}{群概形}的文章中。\footnote{编者注：此处尚待补入评论，包括指向 VI\textsubscript{B} 的参照……}
\begin{SGAresult}{命题}{C.2.1}
设 $S$ 为\Term{discrete-valuation-ring}{离散赋值环}的谱，$t$ 为其\Term{generic-point}{一般点}，$G$ 为在 $S$ 上\Term{finite-type}{有限型}且\Term{flat}{平坦}的 $S$-\Term{group-prescheme}{群预概形}，$\widetilde G=\operatorname{Spec}\Gamma(G,\mathscr O_G)$，$u:G\to\widetilde G$ 为\Term{canonical}{典范}\Term{morphism}{态射}。则：
(1) $\widetilde G$ 自然为 $S$-\Term{group-scheme}{群概形}，$u$ 为\Term{homomorphism}{同态}。
(2) $\operatorname{Ker}(u)$ 在 $S$ 上\Term{flat}{平坦}，且 $(\widetilde G,u)$ 为 $G$ 对 $\operatorname{Ker}(u)$ 的 fpqc 商，故 $\widetilde G$ 是 $G$ 的最大\Term{affine}{仿射}商。
(3) 若 $G_t$ \Term{affine}{仿射}，则 $\operatorname{Ker}(u)$ 是 $S$ 上的\Term{etale}{平展}群，它等于\Term{unit-group}{单位群}当且仅当 $G$ 在 $S$ 上\Term{separated}{分离}。特别地，在 $S$ 上\Term{flat}{平坦}、\Term{finite-type}{有限型}、\Term{separated}{分离}且\Term{generic-fiber}{一般纤维}\Term{affine}{仿射}的 $S$-\Term{group-scheme}{群概形} $G$ 是\Term{affine}{仿射}的。
\end{SGAresult}

\textbf{C.3。}在 XV 6.6 的陈述中，为使满足 $H_s$ 为 $G_s$ 的\Term{parabolic-subgroup}{抛物子群}的 $s\in S$ 构成\Term{open-subset}{开集}，无需假定 $H$ 等于其\Term{connected}{连通}\Term{normalizer}{正规化子}。事实上，回到\Term{lemma}{引理} 6.9 前 c) 段的证明，以 $\overline N$ 表示 $N_t$ 在 $G$ 中的\Term{scheme-theoretic-closure}{概形闭包}。于是 $\overline N$ 为 $G$ 的闭子\Term{group-scheme}{群概形}，在 $S$ 上\Term{flat}{平坦}，包含 $H$ 且包含于 $N$。由下述\Term{theorem}{定理}，$G/\overline N$ \Term{representable}{可表示}。因 $G_s/\overline N_s$ \Term{proper}{固有}且\Term{connected}{连通}，照所引处完成证明即可。
\begin{SGAresult}{\Term{theorem}{定理}}{C.3.1}
\footnote{编者注：参见 S. Anantharaman，《Schémas en groupes, espaces homogènes et espaces algébriques sur une base de dimension 1》，Bull. Soc. Math. France，Mém. 33 (1976)，5--79 中的\Term{theorem}{定理} 4C。}设 $S$ 为 Hensel \Term{discrete-valuation-ring}{离散赋值环}的谱，$G$ 为\Term{locally-finite-type}{局部有限型}的 $S$-\Term{group-prescheme}{群预概形}，$H$ 为 $G$ 的闭子\Term{group-prescheme}{群预概形}，且在 $S$ 上\Term{flat}{平坦}，则 $G/H$ \Term{representable}{可表示}。
\end{SGAresult}

\textbf{C.4. 对 XVI 1.1 的补充。}下述\Term{theorem}{定理}使 VIII 7.9 更为精确。
\begin{SGAresult}{\Term{theorem}{定理}}{C.4.1}
设 $S$ 为\Term{residue-characteristic}{剩余特征} $p>0$ 的\Term{discrete-valuation-ring}{离散赋值环}的谱，$G$ 为在 $S$ 上交换、\Term{smooth}{光滑}、\Term{finite-type}{有限型}且\Term{separated}{分离}的 $S$-\Term{group-scheme}{群概形}。则下列条件\Term{equivalence}{等价}：
(i) ${}_pG$ 在 $S$ 上\Term{finite}{有限}。
(ii) 对每个 $S$-\Term{prescheme}{预概形} $S'$，以及每个在 $S'$ 上\Term{finitely-presented}{有限表现}、\Term{separated}{分离}的 $S'$-\Term{group-prescheme}{群预概形} $H'$，任意 $S'$-\Term{monomorphism}{单态射} $u:G_{S'}\to H'$ 都是\Term{immersion}{浸入}。
\end{SGAresult}
\SourcePage{394}
\textbf{C.5。}例 XVII 6.4 a) 给出了域 $k$ 上的\Term{smooth}{光滑}群，其\Term{unipotent-radical}{幺幂根基}未定义在 $k$ 上。以下命题给出获得这种群的一般方法。
\begin{SGAresult}{命题}{C.5.1}
设 $k$ 为域，$K$ 为 $k$ 的\Term{finite}{有限}\Term{radicial}{根性}\Term{extension}{扩张}，\Term{degree}{次数}大于 $1$；$H$ 为\Term{connected}{连通}、\Term{smooth}{光滑}、\Term{dimension}{维数}为 $r$ 的 $K$-\Term{algebraic-group}{代数群}，$G=\prod_{K/k}H/K$，则 $G$ 为\Term{smooth}{光滑}、\Term{connected}{连通}、\Term{dimension}{维数}为 $r[K:k]$ 的 $k$-\Term{algebraic-group}{代数群}。设 $u:G_K\to H$ 为\Term{canonical}{典范}\Term{homomorphism}{同态}，$R=\operatorname{Ker}(u)$。则：
i) \Term{morphism}{态射} $u$ 为\Term{epimorphism}{满态射}，$R$ 为\Term{smooth}{光滑}、\Term{unipotent}{幺幂}、\Term{connected}{连通}的\Term{algebraic-group}{代数群}。
ii) 若 $U$ 为 $G$ 的\Term{smooth}{光滑}\Term{algebraic-subgroup}{代数子群}，且 $U_K$ 包含 $R$，则 $U=G$。
\end{SGAresult}
\begin{SGAresult}{推论}{}
沿用上述命题的记号。
a) 若 $H$ 不\Term{unipotent}{幺幂}，则 $G_{\bar k}$ 的\Term{unipotent-radical}{幺幂根基}未定义在 $k$ 上。
b) 若 $H$ 为非零\Term{abelian-variety}{阿贝尔簇}，则 $G$ 不是以\Term{smooth}{光滑}\Term{linear-group}{线性群}为核的\Term{abelian-variety}{阿贝尔簇}的\Term{extension}{扩张}。
c) 若 $H$ 不\Term{solvable}{可解}，则 $G_{\bar k}$ 的\Term{solvable-radical}{可解根基}未定义在 $k$ 上，且 $G$ 没有定义在 $k$ 上的 \Term{borel-subgroup}{Borel 子群}。
\end{SGAresult}
先说明如何由命题\Term{pushout}{推出}此推论。
a) 设 $U$ 为 $G$ 的\Term{unipotent-radical}{幺幂根基}。则 $U_K$ 为 $G_K$ 的\Term{unipotent-radical}{幺幂根基}，故包含 $R$，因为由 i)，$R$ \Term{smooth}{光滑}、\Term{unipotent}{幺幂}、\Term{connected}{连通}；于是由 ii)，$U=G$。但 $G_K$ 有商 $H$，按假设 $H$ 不\Term{unipotent}{幺幂}，故 $G$ 不\Term{unipotent}{幺幂}，矛盾。
b) 若 $G$ 是以在 $k$ 上\Term{smooth}{光滑}的\Term{linear-group}{线性群} $L$ 为核的\Term{abelian-variety}{阿贝尔簇}的\Term{extension}{扩张}，则必有 $L_K\supset R$，故 $L=G$。但 $G$ 不可能为\Term{linear-group}{线性群}，因为 $G_K$ 有一个为非零\Term{abelian-variety}{阿贝尔簇}的商。
c) 设 $S$（相应地，$B$）为 $G$ 的\Term{solvable-radical}{可解根基}（相应地，\Term{borel-subgroup}{Borel 子群}），则 $S_K$（相应地，$B_K$）包含 $R$，故 $S=G$（相应地，$B=G$）。但这样 $G$ \Term{solvable}{可解}，与 $G_K$ 有不\Term{solvable}{可解}商矛盾。
\begin{proof}[命题 5.1 的证明]
i) 先描述\Term{morphism}{态射} $u$。给定 $K$-\Term{scheme}{概形} $f:S\to\operatorname{Spec}(K)$，便\Term{constructible}{可构造}\Term{commutative-diagram}{交换图}
\[
\begin{tikzcd}[row sep=large,column sep=large]
\operatorname{Spec}(K)\arrow[d,"j"']&T\arrow[l,"h"']\arrow[d,"j_T",bend right=15]\\
\operatorname{Spec}(k)&S\arrow[l,"g"]\arrow[ul,"f"]\arrow[u,"s",bend right=15]
\end{tikzcd}
\]
\SourcePage{395}
其中 $g=j\circ f$，$T=\operatorname{Spec}(K)\times_{\operatorname{Spec}(k)}S$，$h$ 和 $j_T$ 为两个\Term{projection}{投影}，$s$ 为 $T$ 在 $S$ 上的\Term{section}{截面}，满足 $h\circ s=f$。映射 $u(S):G_K(S)\to H(S)$ 就是复合
\[G_K(S)\xrightarrow{\sim}H(T)\longrightarrow H(S),\]
其中最后一个\Term{arrow}{箭头}由\Term{section}{截面} $s$ 定义。
特别地，取 $S$ 为 $k$ 的\Term{algebraically-closed}{代数闭}包 $\bar k$ 的谱，$f$ 为唯一的 $k$-\Term{morphism}{态射} $K\to\bar k$，于是 $T$ 为\Term{artinian-local-scheme}{局部阿廷概形}。为证明 i)，只须在\Term{base-field}{基域}\Term{extension}{扩张} $K\to\bar k$ 后证明。显然 $G_S=\prod_{T/S}H_T/T$ 表示 $H_T$ 相对于 $S$ 的 Greenberg \Term{functor}{函子}（M. J. Greenberg，《Schemata over local rings》，Ann. of Maths. 73 (1961)，624--648）。由以上描述，在此等同下，$u_S$ 为\Term{canonical}{典范}过渡\Term{morphism}{态射}
\[\operatorname{Green}(H_T)\longrightarrow H_S=H_T\times_TS.\]
断言 i) 由 $H$ 在 $K$ 上\Term{smooth}{光滑}，以及 M. J. Greenberg，《Schemata over local rings II》，Ann. of Maths. 78 (1963)，256--266 得出。
ii) 为证明 ii)，可以假定 $k$ \Term{separable}{可分}闭。设 $U$ 为 $G$ 的\Term{smooth}{光滑}\Term{algebraic-subgroup}{代数子群}，$U_K\supset R$，须证明 $U=G$。因 $G$ \Term{connected}{连通}，可以假定 $U$ \Term{connected}{连通}。设 $V$ 为 $H$ 的\Term{smooth}{光滑}\Term{connected}{连通}\Term{algebraic-subgroup}{代数子群} $u(U_K)$，$V'=\prod_{K/k}V/K$，则 $V'$ 为 $G$ 的\Term{smooth}{光滑}\Term{connected}{连通}\Term{algebraic-subgroup}{代数子群}。群 $V'_K$ 为 $G_K$ 的\Term{algebraic-subgroup}{代数子群}，\Term{canonical}{典范}\Term{morphism}{态射} $V'_K\to V$ 就是 $u$ 在 $V'_K$ 上的\Term{restriction}{限制}。按假设，$U_K$ 包含 $R$，更包含 $\operatorname{Ker}(V'_K\to V)$；另一方面，按构造，$U_K$ 在 $H$ 中的像为 $V$。故 $U_K$ 包含 $V'_K$，从而 $U$ 包含 $V'$。另一方面，\Term{canonical}{典范}\Term{isomorphism}{同构}
\[G(k)\hookrightarrow G(K)\xrightarrow{u(K)}H(K)\]
显然将 $U(k)$ 映入 $V(K)=V'(k)$，即 $U(k)\subset V'(k)$。因 $U$ \Term{smooth}{光滑}且 $k$ \Term{separable}{可分}闭，\Term{pushout}{推出} $U\subset V'$，故 $U=V'$。于是有等式
\[([K:k]-1)\dim V=\dim\operatorname{Ker}(V'_K\to V)=\dim R=([K:k]-1)\dim H.\]
因 $K\ne k$，得 $\dim V=\dim H$，从而 $V=H$，最终 $U=V'=G$。
\end{proof}
\SourcePage{396}
% 原 PDF 第 396 页为空白页。
